\documentclass[a4paper,12pt]{article}
\usepackage[T2A]{fontenc}
\usepackage[utf8]{inputenc}
\usepackage[russian]{babel}
\usepackage{graphicx}
\usepackage{listingsutf8}   % UTF-8 в lstlisting (вместо listings)
\usepackage{xcolor}
\usepackage{geometry}
\usepackage{hyperref}
\usepackage{enumitem}
\usepackage{booktabs}
\usepackage{array}
\usepackage{amsmath}

\geometry{margin=2.5cm}

\hypersetup{
    colorlinks=true,
    linkcolor=blue,
    urlcolor=cyan
}

\lstset{
    language=C++,
    inputencoding=utf8,
    extendedchars=true,
    basicstyle=\ttfamily\small,
    keywordstyle=\color{blue}\bfseries,
    commentstyle=\color{gray}\itshape,
    stringstyle=\color{red!70!black},
    showstringspaces=false,
    breaklines=true,
    frame=single,
    numbers=left,
    numberstyle=\tiny\color{gray},
    backgroundcolor=\color{gray!8},
    tabsize=4,
    captionpos=b,
    xleftmargin=1.5em,
    framexleftmargin=1.5em
}

\title{\textbf{Лекции по C++}\\\large 2 семестр}
\author{Дмитрий Величко}
\date{Март 2026}

\begin{document}
\maketitle
\tableofcontents
\newpage

% -------------------------------------------------------
\section{Исключения в C++}

% -------------------------------------------------------
\subsection{Ключевое слово \texttt{throw}}

\subsubsection{Базовый синтаксис}

Рассмотрим функцию деления:

\begin{lstlisting}
#include <stdexcept>

int divide(int a, int b) {
    if (b == 0) {
        throw std::invalid_argument("division by zero");
    }
    return a / b;
}
\end{lstlisting}

При делении на ноль поведение программы неопределено, поэтому используется
ключевое слово \texttt{throw}. Бросать можно \textbf{любое выражение},
однако принято бросать объекты стандартных исключений из
\texttt{<stdexcept>}.

\texttt{throw} --- один из способов выйти из функции наряду с
\texttt{return}.

\subsubsection{Что можно бросать?}

\textbf{Любой копируемый или перемещаемый объект:}

\begin{lstlisting}
throw 42;                                  // int
throw std::string("error");               // std::string
throw std::runtime_error("fail");         // стандартное исключение
throw 3.14;                               // double
throw std::vector<int>{1, 2, 3};          // даже вектор (но не стоит)

struct MyError { int code; std::string msg; };
throw MyError{404, "not found"};          // пользовательская структура
\end{lstlisting}

\textbf{Рекомендация:} бросайте только наследников \texttt{std::exception}.
Это позволяет единообразно ловить все ошибки через
\texttt{catch(const std::exception\&)}.

\subsubsection{Что происходит при \texttt{throw}?}

\begin{enumerate}
    \item Объект исключения \textbf{копируется} (или перемещается) в специальную
          область памяти, управляемую рантаймом.
    \item Начинается \textbf{раскрутка стека} (stack unwinding).
    \item Управление передаётся первому подходящему блоку \texttt{catch}.
    \item Если подходящего \texttt{catch} нет --- вызывается
          \texttt{std::terminate()}.
\end{enumerate}

\subsubsection{\texttt{throw} как выражение}

\texttt{throw} --- это выражение типа \texttt{void}. Его можно использовать
в тернарном операторе и других контекстах:

\begin{lstlisting}
int safe_divide(int a, int b) {
    return b != 0 ? a / b : throw std::invalid_argument("zero divisor");
}

// В лямбде:
auto get_or_throw = [](std::optional<int> opt) -> int {
    return opt.has_value() ? *opt : throw std::runtime_error("empty");
};
\end{lstlisting}

% -------------------------------------------------------
\subsection{Раскрутка стека (Stack Unwinding)}

\subsubsection{Определение}

При выбросе исключения происходит \textbf{раскрутка стека} (stack
unwinding) --- автоматическое уничтожение всех локальных объектов
в обратном порядке их создания, с вызовом деструкторов, по цепочке
вызовов от точки \texttt{throw} до подходящего блока \texttt{catch}.

\subsubsection{Пошаговая механика}

\begin{lstlisting}
#include <iostream>
#include <stdexcept>

struct Guard {
    std::string name;
    Guard(std::string n) : name(std::move(n)) {
        std::cout << "  + " << name << "\n";
    }
    ~Guard() {
        std::cout << "  - " << name << "\n";
    }
};

void f3() {
    Guard g("f3_local");
    throw std::runtime_error("boom");  // (1) Бросаем исключение
    // Код ниже НИКОГДА не выполнится
    std::cout << "after throw\n";
}

void f2() {
    Guard g("f2_local");
    f3();                               // (2) Вызываем f3
    std::cout << "after f3\n";          // Не выполнится
}

void f1() {
    Guard g("f1_local");
    f2();                               // (3) Вызываем f2
    std::cout << "after f2\n";          // Не выполнится
}

int main() {
    try {
        f1();
    } catch (const std::exception& e) {
        std::cout << "Caught: " << e.what() << "\n";
    }
}
\end{lstlisting}

\textbf{Вывод:}
\begin{lstlisting}[numbers=none]
  + f1_local
  + f2_local
  + f3_local
  - f3_local       <- раскрутка: уничтожается f3_local
  - f2_local       <- раскрутка: уничтожается f2_local
  - f1_local       <- раскрутка: уничтожается f1_local
Caught: boom
\end{lstlisting}

Обратите внимание: деструкторы вызываются в \textbf{обратном порядке} создания
(LIFO), а строки \texttt{"after ..."} не выполняются --- управление <<перепрыгивает>>
через них.

\subsubsection{Визуализация стека}

\begin{lstlisting}[numbers=none]
Стек вызовов в момент throw:

|  f3()  | <- throw runtime_error("boom")
|  f2()  |    ~Guard("f3_local")
|  f1()  |    ~Guard("f2_local")
| main() |    ~Guard("f1_local")
+--------+    catch (const std::exception& e) <- сюда попадаем
\end{lstlisting}

Рантайм <<разматывает>> стек снизу вверх: для каждого кадра уничтожает
локальные объекты, затем переходит к следующему, пока не найдёт \texttt{catch}.

\subsubsection{Что уничтожается, а что нет}

\begin{center}
\begin{tabular}{lcc}
\toprule
\textbf{Объект} & \textbf{Уничтожается?} & \textbf{Почему} \\
\midrule
Локальные переменные на стеке  & Да  & RAII \\
Объекты, созданные \texttt{new}  & Нет & Нет владельца на стеке \\
\texttt{std::unique\_ptr}        & Да  & Деструктор вызовет \texttt{delete} \\
\texttt{std::shared\_ptr}        & Да  & Деструктор уменьшит счётчик \\
Сырые указатели (\texttt{int*})  & Да\textsuperscript{*} & Но память \textbf{не} освобождается! \\
\bottomrule
\end{tabular}
\end{center}

\textsuperscript{*}Переменная-указатель уничтожается, но \texttt{delete} для неё
никто не вызывает --- это \textbf{утечка памяти}.

\begin{lstlisting}
void leak_example() {
    int* raw = new int(42);          // Выделили память
    std::unique_ptr<int> safe(new int(43));  // RAII-обёртка

    throw std::runtime_error("oops");

    delete raw;  // НИКОГДА не выполнится -> утечка!
    // safe -- деструктор unique_ptr вызовет delete автоматически
}
\end{lstlisting}

\textbf{Мораль:} именно поэтому RAII критически важен. Любой ресурс
(память, файл, мьютекс, соединение) должен быть обёрнут в объект с деструктором.

\subsubsection{Частично сконструированные объекты}

Если исключение вылетает \textbf{из конструктора}, деструктор этого объекта
\textbf{не вызывается} (объект не считается созданным). Но деструкторы уже
инициализированных полей \textbf{вызываются}:

\begin{lstlisting}
struct Resource {
    Resource(int id) { std::cout << "R(" << id << ")\n"; }
    ~Resource()      { std::cout << "~R\n"; }
};

struct Widget {
    Resource a, b, c;

    Widget() : a(1), b(2), c(3) {
        // Допустим, конструктор c бросил исключение
        // Тогда:
        //   a -- уничтожается (был создан)
        //   b -- уничтожается (был создан)
        //   c -- НЕ уничтожается (не был до конца создан)
        //   ~Widget() -- НЕ вызывается (Widget не был создан)
    }
};
\end{lstlisting}

\begin{lstlisting}
struct Dangerous {
    Resource a{1};
    Resource b{2};

    Dangerous() {
        Resource* raw = new Resource(3);
        throw std::runtime_error("oops");
        delete raw;  // Утечка! ~Dangerous() не вызовется,
                     // но ~a и ~b вызовутся
    }
};
\end{lstlisting}

\subsubsection{Раскрутка стека и массивы}

При раскрутке массива деструкторы вызываются только для
\textbf{уже сконструированных} элементов:

\begin{lstlisting}
struct Item {
    static int count;
    Item() {
        if (++count == 4) throw std::runtime_error("4th item");
        std::cout << "Item(" << count << ")\n";
    }
    ~Item() { std::cout << "~Item\n"; }
};
int Item::count = 0;

void test() {
    Item arr[5];
    // Создаётся Item(1), Item(2), Item(3)
    // Item(4) бросает исключение
    // Раскрутка: ~Item для arr[2], arr[1], arr[0] (в обратном порядке)
    // arr[3] и arr[4] не были созданы -- деструкторы не вызываются
}
\end{lstlisting}

\subsubsection{Когда раскрутка НЕ происходит}

\begin{itemize}
    \item \textbf{Не пойманное исключение:} если нигде нет подходящего
          \texttt{catch}, стандарт гарантирует только вызов
          \texttt{std::terminate()}. Раскрутка стека
          \textit{implementation-defined} --- может произойти, а может и нет.
    \item \textbf{\texttt{std::abort()} / \texttt{std::\_Exit()}:} немедленное
          завершение без раскрутки.
    \item \textbf{Нарушение \texttt{noexcept}:} если функция помечена
          \texttt{noexcept} и из неё вылетает исключение ---
          \texttt{std::terminate()} без гарантированной раскрутки.
\end{itemize}

% -------------------------------------------------------
\subsection{try / catch}

\subsubsection{Базовый синтаксис}

\begin{lstlisting}
#include <iostream>
#include <stdexcept>

int main() {
    try {
        divide(1, 0);
    } catch (const std::exception& e) {
        std::cout << e.what() << "\n";
    }
}
\end{lstlisting}

\begin{itemize}
    \item Метод \texttt{what()} определён в \texttt{std::exception}
          и возвращает строку с описанием ошибки --- ту самую, что
          была передана при создании объекта исключения.
    \item Исключение следует ловить \textbf{по константной ссылке},
          чтобы избежать лишнего копирования и сохранить полиморфизм.
\end{itemize}

\subsubsection{Множественные блоки \texttt{catch}}

\begin{lstlisting}
try {
    do_something();
} catch (const std::out_of_range& e) {
    // Обработка выхода за границы
    std::cerr << "Index error: " << e.what() << "\n";
} catch (const std::invalid_argument& e) {
    // Обработка неверного аргумента
    std::cerr << "Bad argument: " << e.what() << "\n";
} catch (const std::exception& e) {
    // Все остальные стандартные исключения
    std::cerr << "Error: " << e.what() << "\n";
} catch (...) {
    // Абсолютно любое исключение (включая int, string и т.д.)
    std::cerr << "Unknown error\n";
}
\end{lstlisting}

\textbf{Правила:}
\begin{itemize}
    \item Блоки \texttt{catch} проверяются \textbf{сверху вниз},
          побеждает \textbf{первый подходящий}.
    \item \texttt{catch(...)} ловит \textbf{всё}, поэтому должен быть последним.
    \item Более специфичные обработчики ставятся \textbf{перед} более общими.
\end{itemize}

\subsubsection{Ловля по значению vs по ссылке}

\begin{lstlisting}
// ПЛОХО: ловля по значению -- срезка (slicing)!
try {
    throw std::runtime_error("specific error");
} catch (std::exception e) {
    // e -- копия, тип СРЕЗАН до std::exception
    // e.what() может вернуть не то, что ожидали
}

// ХОРОШО: ловля по константной ссылке
try {
    throw std::runtime_error("specific error");
} catch (const std::exception& e) {
    // e -- ссылка на исходный объект, полиморфизм сохранён
    // e.what() вернёт "specific error"
}
\end{lstlisting}

При ловле по значению происходит \textbf{object slicing}: если бросили
наследника, а поймали базовый тип, часть данных наследника \textbf{теряется}.

\subsubsection{Вложенные try/catch}

\begin{lstlisting}
void process() {
    try {
        try {
            throw std::runtime_error("inner");
        } catch (const std::logic_error& e) {
            std::cout << "logic: " << e.what() << "\n";  // Не подходит
        }
        // runtime_error не пойман внутренним catch -- летит дальше
    } catch (const std::exception& e) {
        std::cout << "outer: " << e.what() << "\n";  // "outer: inner"
    }
}
\end{lstlisting}

% -------------------------------------------------------
\subsection{Приведение типов в \texttt{catch}}

В \texttt{catch} \textbf{не работают} обычные неявные преобразования
и правила перегрузки. Доступны только:
\begin{itemize}
    \item Добавление \texttt{const} / приведение к ссылке
    \item Приведение типа потомка к типу родителя
\end{itemize}

\begin{lstlisting}
try {
    throw 1;  // int
} catch (double d) {
    std::cout << "double\n";    // не поймает: int -> double не работает
} catch (long long l) {
    std::cout << "long long\n"; // не поймает: int -> long long не работает
} catch (...) {
    std::cout << "other\n";     // поймает
}
\end{lstlisting}

\subsubsection{Наследование в \texttt{catch}}

Блоки \texttt{catch} проверяются \textbf{сверху вниз}, побеждает
первый подходящий:

\begin{lstlisting}
struct Mom {};
struct Son : Mom {};

int main() {
    try {
        throw Son{};
    } catch (const Mom& m) {
        std::cout << "Mom\n"; // поймает: Son приводится к Mom
    } catch (const Son& s) {
        std::cout << "Son\n"; // сюда не дойдём
    } catch (...) {
        std::cout << "other\n";
    }
}
\end{lstlisting}

\textbf{Вывод:} \texttt{Mom}. Именно поэтому обработчики потомков
нужно ставить \textbf{перед} обработчиками родителей, иначе они
никогда не сработают.

\subsubsection{Правильный порядок}

\begin{lstlisting}
try {
    throw Son{};
} catch (const Son& s) {
    std::cout << "Son\n";  // Теперь поймает!
} catch (const Mom& m) {
    std::cout << "Mom\n";  // Ловит остальных наследников Mom
} catch (...) {
    std::cout << "other\n";
}
\end{lstlisting}

% -------------------------------------------------------
\subsection{Где встречаются исключения в стандартной библиотеке}

\begin{center}
\begin{tabular}{lp{7cm}}
\toprule
\textbf{Операция} & \textbf{Исключение} \\
\midrule
\texttt{vector::at(i)}
    & \texttt{std::out\_of\_range} при \texttt{i >= size()} \\
\addlinespace
\texttt{operator new}
    & \texttt{std::bad\_alloc} при нехватке памяти \\
\addlinespace
\texttt{dynamic\_cast<T\&>(ref)}
    & \texttt{std::bad\_cast} при неудачном приведении \\
\addlinespace
\texttt{std::stoi("abc")}
    & \texttt{std::invalid\_argument} (не число) \\
\addlinespace
\texttt{std::stoi("99999999999")}
    & \texttt{std::out\_of\_range} (переполнение) \\
\addlinespace
\texttt{map::at(key)}
    & \texttt{std::out\_of\_range} при отсутствии ключа \\
\addlinespace
\texttt{optional::value()}
    & \texttt{std::bad\_optional\_access} если пуст \\
\addlinespace
\texttt{variant::get<T>()}
    & \texttt{std::bad\_variant\_access} при неверном типе \\
\addlinespace
\texttt{thread::join()} на отсоединённом
    & \texttt{std::system\_error} \\
\bottomrule
\end{tabular}
\end{center}

% -------------------------------------------------------
\subsection{Runtime-ошибки vs Исключения C++}

Это \textbf{принципиально разные вещи}.

\subsubsection{Три вида runtime-ошибок на уровне ОС}

\begin{center}
\begin{tabular}{lp{9cm}}
\toprule
\textbf{Сигнал} & \textbf{Типичные причины} \\
\midrule
\texttt{SIGSEGV} (Segfault)
    & Выход за пределы массива, разыменование \texttt{nullptr},
      переполнение стека \\
\addlinespace
\texttt{SIGFPE} (Floating Point Exception)
    & Деление целого на ноль, переполнение целых \\
\addlinespace
\texttt{SIGABRT} (Aborted)
    & Вызов \texttt{std::abort()}, \texttt{std::terminate()},
      падение \texttt{assert} \\
\bottomrule
\end{tabular}
\end{center}

\medskip
Эти ошибки \textbf{не являются} исключениями C++ и \textbf{не ловятся}
через \texttt{catch}:

\begin{lstlisting}
int arr[5];
try {
    arr[10] = 42;  // SIGSEGV -- catch не поймает
} catch (...) {
    // сюда мы НЕ попадём
}

try {
    int x = 1 / 0;  // SIGFPE -- тоже не исключение C++
} catch (...) {
    // и сюда не попадём
}
\end{lstlisting}

\textbf{Важно:} деление на ноль --- это исключение на уровне
\textit{процессора}, не языка. Процессор обращается к ОС, которая
посылает \texttt{SIGFPE}. Это совпадение терминологии, но разные
уровни абстракции.

\subsubsection{Пути к \texttt{std::terminate()}}

\begin{itemize}
    \item Не пойманное исключение (uncaught exception)
    \item Явный вызов \texttt{std::terminate()} или \texttt{std::abort()}
    \item Упавший \texttt{assert} (он напрямую вызывает \texttt{abort()})
    \item Исключение из \texttt{noexcept}-функции
    \item Исключение из деструктора во время раскрутки стека
    \item Два одновременных активных исключения (double exception)
\end{itemize}

С точки зрения ОС программа, завершившаяся из-за не пойманного C++-исключения,
получает вердикт \texttt{SIGABRT}.

Также стоит отметить, что \texttt{std::runtime\_error} ---
это \textit{частный случай} \texttt{std::exception}, то есть обычное
C++-исключение, не имеющее отношения к SIGFPE/SIGSEGV/SIGABRT.

% -------------------------------------------------------
\subsection{Копирование при броске исключения}

\begin{lstlisting}
#include <iostream>

struct A {
    A()         { std::cout << "A()\n"; }
    ~A()        { std::cout << "~A()\n"; }
    A(const A&) { std::cout << "A(copy)\n"; }
};

void f(int x) {
    A a;
    if (x == 0) {
        throw a;  // объект копируется в специальную память
    }
}

int main() {
    try {
        f(0);
    } catch (...) {
        std::cout << "caught!\n";
    }
}
\end{lstlisting}

\textbf{Порядок вывода:}
\begin{enumerate}
    \item \texttt{A()} --- создание \texttt{a} на стеке
    \item \texttt{A(copy)} --- копирование в специальную память при \texttt{throw}
    \item \texttt{\textasciitilde A()} --- уничтожение \texttt{a} при раскрутке стека
    \item \texttt{caught!}
    \item \texttt{\textasciitilde A()} --- уничтожение копии при выходе из \texttt{catch}
\end{enumerate}

\medskip
\begin{itemize}
    \item \texttt{catch(A a)} --- ещё одно копирование из дин.\ памяти на стек.
    \item \texttt{catch(A\& a)} --- работаем прямо с объектом исключения, без копирования.
    \item \texttt{catch(const A\& a)} --- то же, но без возможности модификации
          (рекомендуемый вариант).
\end{itemize}

\subsubsection{Где хранится объект исключения?}

Стандарт не специфицирует точное место хранения. На практике ---
в \textbf{динамической памяти}. Исключение: \texttt{std::bad\_alloc} ---
при нехватке памяти нельзя выделить новую, поэтому компилятор заранее
резервирует буфер в \textbf{статической памяти}
(\textit{emergency buffer}).

Исключения дороги по стоимости: выделение динамической памяти происходит
в любом случае. Поэтому их не стоит использовать в горячих путях кода.

\subsubsection{Оптимизация: \texttt{throw} временного объекта}

\begin{lstlisting}
// Лишнее копирование:
void bad() {
    std::runtime_error err("fail");
    throw err;  // Копирование err в специальную память
}

// Без лишнего копирования:
void good() {
    throw std::runtime_error("fail");  // Объект создаётся сразу в спец. памяти
                                        // Компилятор может применить copy elision
}
\end{lstlisting}

\subsubsection{Повторный проброс}

\begin{lstlisting}
try {
    f(0);
} catch (const A& a) {
    throw;      // прокидываем то же исключение дальше (без копирования)
    // throw a; // создаём НОВЫЙ объект -- возможна срезка типа!
}
\end{lstlisting}

\textbf{Критическая разница:}
\begin{itemize}
    \item \texttt{throw;} --- пробрасывает \textbf{оригинальный} объект исключения.
          Тип сохраняется. Нет копирования.
    \item \texttt{throw a;} --- создаёт \textbf{новый} объект, копируя \texttt{a}.
          Если реальный тип --- наследник, происходит \textbf{object slicing}!
\end{itemize}

\begin{lstlisting}
struct Base { virtual const char* what() const { return "Base"; } };
struct Derived : Base { const char* what() const override { return "Derived"; } };

void demonstrate_slicing() {
    try {
        throw Derived{};
    } catch (const Base& e) {
        std::cout << e.what() << "\n";  // "Derived" (полиморфизм)
        // throw;                        // Пробросит Derived -- правильно
        throw e;                         // Пробросит Base -- СРЕЗКА!
    }
}
\end{lstlisting}

\textbf{Правило: всегда используйте \texttt{throw;} без аргумента для повторного
проброса.}

Если исключение не поймано нигде --- гарантируется только вызов
\texttt{std::terminate()}. Деструкторы при этом могут не вызваться.

% -------------------------------------------------------
\subsection{Спецификатор \texttt{noexcept}}

\subsubsection{Что такое \texttt{noexcept}?}

\texttt{noexcept} --- спецификатор, обещающий компилятору, что функция
\textbf{не бросает исключений}:

\begin{lstlisting}
void safe_function() noexcept {
    // Если отсюда вылетит исключение -- std::terminate()
}

int add(int a, int b) noexcept {
    return a + b;  // Действительно не бросает
}
\end{lstlisting}

\subsubsection{Зачем это нужно?}

\begin{enumerate}
    \item \textbf{Оптимизация:} компилятор может генерировать более
          эффективный код, зная, что исключений не будет (не нужна таблица
          раскрутки для этой функции).
    \item \textbf{Семантика \texttt{std::vector}:} при реаллокации вектор
          использует \texttt{std::move\_if\_noexcept}. Если конструктор перемещения
          помечен \texttt{noexcept} --- элементы перемещаются. Если нет ---
          \textbf{копируются} (для strong exception safety).
    \item \textbf{Документация:} явное указание контракта функции.
\end{enumerate}

\begin{lstlisting}
class MyString {
public:
    // Без noexcept -- vector будет КОПИРОВАТЬ при реаллокации
    MyString(MyString&& other) { /* ... */ }

    // С noexcept -- vector будет ПЕРЕМЕЩАТЬ (быстрее!)
    // MyString(MyString&& other) noexcept { /* ... */ }
};

std::vector<MyString> vec;
vec.push_back(MyString());  // Реаллокация: move или copy?
\end{lstlisting}

\subsubsection{Условный \texttt{noexcept}}

\texttt{noexcept} может принимать выражение:

\begin{lstlisting}
template<typename T>
void swap_values(T& a, T& b) noexcept(noexcept(T(std::move(a)))) {
    T tmp = std::move(a);
    a = std::move(b);
    b = std::move(tmp);
}
// noexcept, если конструктор перемещения T -- noexcept
\end{lstlisting}

\texttt{noexcept(expr)} как оператор возвращает \texttt{true}/\texttt{false}:

\begin{lstlisting}
static_assert(noexcept(42 + 1));           // true -- арифметика не бросает
static_assert(!noexcept(std::stoi("42"))); // false -- stoi может бросить

void f() noexcept;
static_assert(noexcept(f()));              // true -- f помечена noexcept
\end{lstlisting}

\subsubsection{Что будет, если \texttt{noexcept}-функция бросит?}

\begin{lstlisting}
void oops() noexcept {
    throw std::runtime_error("surprise!");  // Компилируется!
}

int main() {
    try {
        oops();
    } catch (...) {
        // Сюда мы НЕ попадём!
    }
    // Программа завершится через std::terminate()
}
\end{lstlisting}

Компилятор \textbf{не проверяет} тело \texttt{noexcept}-функции на наличие
\texttt{throw}. \texttt{noexcept} --- это \textbf{обещание} программиста,
а не проверка компилятора. Нарушение --- \texttt{std::terminate()}.

\subsubsection{Что помечать \texttt{noexcept}?}

\begin{center}
\begin{tabular}{lc}
\toprule
\textbf{Функция} & \texttt{noexcept}? \\
\midrule
Деструкторы                    & Да (по умолчанию с C++11) \\
Move-конструкторы              & Да, если возможно \\
Move-присваивание              & Да, если возможно \\
\texttt{swap()}                & Да, если возможно \\
Функции, работающие с памятью  & Зависит \\
Функции с внешними вызовами    & Обычно нет \\
\bottomrule
\end{tabular}
\end{center}

\textbf{Важно:} деструкторы в C++11 и новее \textbf{неявно}
\texttt{noexcept}. Исключение из деструктора --- почти всегда
\texttt{std::terminate()}.

% -------------------------------------------------------
\subsection{Исключения в деструкторах}

\subsubsection{Почему нельзя бросать из деструктора?}

Деструкторы вызываются при раскрутке стека. Если во время раскрутки
(уже есть активное исключение) деструктор бросит \textbf{второе} исключение ---
в C++ не может быть двух активных исключений одновременно, и вызывается
\texttt{std::terminate()}:

\begin{lstlisting}
struct Bad {
    ~Bad() noexcept(false) {
        throw std::runtime_error("from destructor");
    }
};

void example() {
    try {
        Bad b;
        throw std::runtime_error("first");
        // Начинается раскрутка стека
        // Вызывается ~Bad(), который бросает ВТОРОЕ исключение
        // -> std::terminate()!
    } catch (...) {
        // Сюда не дойдём
    }
}
\end{lstlisting}

\subsubsection{Как обрабатывать ошибки в деструкторе?}

\begin{lstlisting}
class FileWriter {
    FILE* file_;
public:
    ~FileWriter() {
        // Вариант 1: проглотить исключение (обычный подход)
        try {
            if (file_ && fclose(file_) != 0) {
                // Логируем ошибку, но НЕ бросаем
                std::cerr << "Warning: failed to close file\n";
            }
        } catch (...) {
            // Никогда не позволяем исключению покинуть деструктор
        }
    }

    // Вариант 2: явный метод close() -- пользователь вызовет заранее
    void close() {
        if (file_ && fclose(file_) != 0) {
            throw std::runtime_error("close failed");
        }
        file_ = nullptr;
    }
};
\end{lstlisting}

\subsubsection{Функция \texttt{std::uncaught\_exceptions()}}

C++17 добавил \texttt{std::uncaught\_exceptions()} (с~<<s>>!), которая
возвращает \textbf{количество} активных непойманных исключений. Это позволяет
деструктору узнать, вызван ли он в процессе раскрутки:

\begin{lstlisting}
class ScopeGuard {
    int exceptions_on_entry_;
    std::function<void()> on_success_;
    std::function<void()> on_failure_;
public:
    ScopeGuard(std::function<void()> success, std::function<void()> failure)
        : exceptions_on_entry_(std::uncaught_exceptions())
        , on_success_(std::move(success))
        , on_failure_(std::move(failure))
    {}

    ~ScopeGuard() {
        if (std::uncaught_exceptions() > exceptions_on_entry_) {
            on_failure_();   // Деструктор вызван из-за исключения
        } else {
            on_success_();   // Нормальный выход из области видимости
        }
    }
};

void test() {
    ScopeGuard guard(
        []{ std::cout << "Success!\n"; },
        []{ std::cout << "Failure!\n"; }
    );
    throw std::runtime_error("oops");  // "Failure!"
}
\end{lstlisting}

% -------------------------------------------------------
\subsection{Function-try-block}

Особый синтаксис \texttt{try/catch} для \textbf{всего тела функции},
включая список инициализации конструктора:

\begin{lstlisting}
class Connection {
    DatabaseHandle db_;
    Logger log_;

public:
    Connection(const std::string& url)
    try : db_(url), log_("connection.log")  // Если db_ или log_ бросят --
    {                                        // поймаем ниже
        db_.connect();
    }
    catch (const std::exception& e) {
        std::cerr << "Failed to create Connection: " << e.what() << "\n";
        // ВАЖНО: для конструктора исключение автоматически перебрасывается!
        // Объект НЕ был создан, вернуть его нельзя
    }
};
\end{lstlisting}

\textbf{Особенности function-try-block в конструкторах:}
\begin{itemize}
    \item В \texttt{catch} блоке \textbf{нельзя} обратиться к полям объекта ---
          они уже могут быть уничтожены.
    \item Исключение \textbf{автоматически перебрасывается} при выходе из
          \texttt{catch}. Подавить его невозможно --- объект не был создан.
    \item Можно заменить исключение на другое через \texttt{throw other;}.
\end{itemize}

Для обычных функций function-try-block просто оборачивает всё тело:

\begin{lstlisting}
int parse(const std::string& s)
try {
    return std::stoi(s);
}
catch (const std::exception& e) {
    std::cerr << "Parse failed: " << e.what() << "\n";
    return -1;  // Для обычных функций -- можно вернуть значение
}
\end{lstlisting}

% -------------------------------------------------------
\subsection{Гарантии безопасности исключений (Exception Safety)}

При проектировании функций и классов важно определить, какие \textbf{гарантии}
они предоставляют в случае исключения.

\subsubsection{Четыре уровня гарантий}

\begin{center}
\begin{tabular}{lp{4cm}p{5cm}}
\toprule
\textbf{Гарантия} & \textbf{Суть} & \textbf{Пример} \\
\midrule
\textbf{No-throw}
    & Функция никогда не бросает исключений
    & Деструкторы, \texttt{swap()}, \texttt{noexcept}-функции \\
\addlinespace
\textbf{Strong} (строгая)
    & При исключении состояние объекта откатывается к исходному (<<всё или ничего>>)
    & \texttt{vector::push\_back}, \texttt{std::atomic} операции \\
\addlinespace
\textbf{Basic} (базовая)
    & При исключении объект в \textit{валидном}, но \textit{неопределённом} состоянии. Утечек нет
    & \texttt{std::string::operator=} (если аллокация не удалась) \\
\addlinespace
\textbf{No guarantee}
    & Никаких гарантий. Возможны утечки и повреждение данных
    & Код без RAII \\
\bottomrule
\end{tabular}
\end{center}

\subsubsection{Пример: строгая гарантия в \texttt{push\_back}}

\begin{lstlisting}
template<typename T>
void vector<T>::push_back(const T& value) {
    if (sz_ == cap_) {
        // 1. Выделяем новую память
        T* newarr = allocate(cap_ * 2);  // Может бросить bad_alloc
                                           // Но старые данные не тронуты!

        // 2. Копируем/перемещаем элементы
        try {
            copy_or_move(arr_, newarr, sz_);  // Может бросить
        } catch (...) {
            deallocate(newarr);  // Откатываемся: освобождаем новую память
            throw;               // Пробрасываем -- старый вектор цел!
        }

        // 3. Успех: заменяем старое на новое
        deallocate(arr_);
        arr_ = newarr;
        cap_ *= 2;
    }
    new (arr_ + sz_) T(value);  // Может бросить -- но вектор в валидном состоянии
    ++sz_;
}
\end{lstlisting}

\subsubsection{Copy-and-swap idiom --- строгая гарантия для \texttt{operator=}}

\begin{lstlisting}
class MyString {
    char* data_;
    size_t size_;
public:
    // Строгая гарантия через copy-and-swap:
    MyString& operator=(MyString other) {  // (1) Копия -- может бросить
        swap(other);                        // (2) Swap -- noexcept
        return *this;                       // (3) Старые данные в other
    }                                       //     other уничтожится в деструкторе

    void swap(MyString& other) noexcept {
        std::swap(data_, other.data_);
        std::swap(size_, other.size_);
    }
};
\end{lstlisting}

Идея: копирование (шаг 1) может бросить, но исходный объект ещё не тронут.
\texttt{swap} (шаг 2) гарантированно не бросает. Если шаг 1 бросил ---
исходный объект остался в прежнем состоянии (строгая гарантия).

% -------------------------------------------------------
\subsection{\texttt{std::exception\_ptr} --- исключения между потоками}

\texttt{std::exception\_ptr} позволяет \textbf{захватить} исключение,
сохранить его и \textbf{перебросить} позже --- в том числе в другом потоке.

\subsubsection{Основные функции}

\begin{center}
\begin{tabular}{lp{8cm}}
\toprule
\textbf{Функция} & \textbf{Описание} \\
\midrule
\texttt{std::current\_exception()}
    & Захватывает текущее исключение (внутри \texttt{catch}) и возвращает
      \texttt{exception\_ptr} \\
\addlinespace
\texttt{std::rethrow\_exception(ptr)}
    & Перебрасывает исключение, хранимое в \texttt{ptr} \\
\addlinespace
\texttt{std::make\_exception\_ptr(e)}
    & Создаёт \texttt{exception\_ptr} из объекта \texttt{e} без
      входа в \texttt{catch} \\
\bottomrule
\end{tabular}
\end{center}

\subsubsection{Пример: передача исключения между потоками}

\begin{lstlisting}
#include <exception>
#include <thread>
#include <iostream>

void worker(std::exception_ptr& eptr) {
    try {
        // Какая-то работа, которая может бросить
        throw std::runtime_error("error in worker thread");
    } catch (...) {
        eptr = std::current_exception();  // Захватываем исключение
    }
}

int main() {
    std::exception_ptr eptr = nullptr;

    std::thread t(worker, std::ref(eptr));
    t.join();

    // Теперь в основном потоке:
    if (eptr) {
        try {
            std::rethrow_exception(eptr);  // Перебрасываем
        } catch (const std::exception& e) {
            std::cout << "From worker: " << e.what() << "\n";
            // "From worker: error in worker thread"
        }
    }
}
\end{lstlisting}

\subsubsection{Пример: сбор нескольких исключений}

\begin{lstlisting}
#include <exception>
#include <vector>

std::vector<std::exception_ptr> errors;

void process_item(int item) {
    try {
        if (item % 3 == 0)
            throw std::runtime_error("bad item: " + std::to_string(item));
        // Обычная обработка...
    } catch (...) {
        errors.push_back(std::current_exception());
    }
}

int main() {
    for (int i = 0; i < 10; ++i) {
        process_item(i);
    }

    std::cout << "Errors: " << errors.size() << "\n";
    for (auto& eptr : errors) {
        try {
            std::rethrow_exception(eptr);
        } catch (const std::exception& e) {
            std::cout << "  - " << e.what() << "\n";
        }
    }
}
\end{lstlisting}

% -------------------------------------------------------
\subsection{\texttt{std::nested\_exception} --- вложенные исключения}

Иногда при обработке одного исключения возникает другое. C++11 позволяет
\textbf{вложить} оригинальное исключение внутрь нового:

\begin{lstlisting}
#include <exception>
#include <stdexcept>
#include <iostream>

void low_level() {
    throw std::runtime_error("disk I/O error");
}

void mid_level() {
    try {
        low_level();
    } catch (...) {
        // std::throw_with_nested оборачивает текущее исключение внутрь нового
        std::throw_with_nested(
            std::runtime_error("mid_level failed")
        );
    }
}

void high_level() {
    try {
        mid_level();
    } catch (...) {
        std::throw_with_nested(
            std::runtime_error("high_level failed")
        );
    }
}
\end{lstlisting}

\subsubsection{Распаковка вложенных исключений}

\begin{lstlisting}
void print_exception_chain(const std::exception& e, int depth = 0) {
    std::cout << std::string(depth * 2, ' ') << e.what() << "\n";
    try {
        std::rethrow_if_nested(e);  // Перебрасывает вложенное, если есть
    } catch (const std::exception& nested) {
        print_exception_chain(nested, depth + 1);
    } catch (...) {
        std::cout << std::string((depth + 1) * 2, ' ') << "unknown\n";
    }
}

int main() {
    try {
        high_level();
    } catch (const std::exception& e) {
        print_exception_chain(e);
    }
}
\end{lstlisting}

\textbf{Вывод:}
\begin{lstlisting}[numbers=none]
high_level failed
  mid_level failed
    disk I/O error
\end{lstlisting}

Это аналог <<exception chaining>> в Java и Python.

% -------------------------------------------------------
\subsection{Стоимость исключений: zero-cost exception model}

Современные компиляторы (GCC, Clang) используют \textbf{table-based}
(zero-cost) модель исключений:

\begin{center}
\begin{tabular}{lp{5cm}p{5cm}}
\toprule
& \textbf{Нормальный путь} (без исключений) & \textbf{При броске исключения} \\
\midrule
Накладные расходы
    & \textbf{Нулевые}. Никаких дополнительных инструкций
    & Высокие: поиск по таблицам, раскрутка стека \\
\addlinespace
Размер бинарника
    & Увеличивается (таблицы раскрутки хранятся в секции \texttt{.eh\_frame})
    & --- \\
\bottomrule
\end{tabular}
\end{center}

\textbf{Принцип:} <<платите только когда бросаете>>. На нормальном пути
исполнения нет вообще никаких затрат --- ни проверок, ни сохранения контекста.
Но когда исключение летит, стоимость \textbf{очень высокая} (в 10--100 раз
дороже обычного \texttt{return}).

\begin{lstlisting}[numbers=none]
Бросок исключения (приблизительно):
  1. Выделение памяти для объекта исключения        ~100 нс
  2. Поиск обработчика по таблицам (DWARF unwind)   ~1-10 мкс
  3. Раскрутка стека (вызов деструкторов)            ~1-10 мкс
  Итого: ~2-20 мкс

Обычный return:                                     ~1 нс
\end{lstlisting}

\textbf{Вывод:} исключения --- для \textbf{исключительных} ситуаций,
а не для управления потоком выполнения. Для ожидаемых ошибок (парсинг,
валидация) используйте \texttt{std::expected} или коды возврата.

% -------------------------------------------------------
\subsection{Итого: чек-лист по исключениям}

\begin{center}
\begin{tabular}{lp{9cm}}
\toprule
\textbf{Правило} & \textbf{Почему} \\
\midrule
Бросайте наследников \texttt{std::exception}
    & Единообразная обработка \\
\addlinespace
Ловите по \texttt{const\&}
    & Нет копирования, нет срезки \\
\addlinespace
Пробрасывайте через \texttt{throw;} (без аргумента)
    & Сохранение типа, нет срезки \\
\addlinespace
Помечайте деструкторы \texttt{noexcept}
    & Они и так \texttt{noexcept} с C++11, но не бросайте из них \\
\addlinespace
Помечайте move-конструкторы \texttt{noexcept}
    & Иначе \texttt{vector} будет копировать, а не перемещать \\
\addlinespace
Используйте RAII для ресурсов
    & Раскрутка стека вызывает только деструкторы \\
\addlinespace
Не используйте исключения в горячих путях
    & Бросок в $\sim$1000$\times$ дороже \texttt{return} \\
\addlinespace
Обработчики потомков --- перед родителями
    & \texttt{catch} проверяется сверху вниз \\
\bottomrule
\end{tabular}
\end{center}

% -------------------------------------------------------
\newpage
\section{Контейнеры и итераторы}

\subsection{\texttt{std::vector}: реализация изнутри}

Начнём с устройства \texttt{std::vector} и в первую очередь рассмотрим
реализацию метода \texttt{push\_back}.

Структура вектора состоит из трёх полей: указателя на данные \texttt{arr\_},
размера \texttt{sz\_} и ёмкости \texttt{cap\_}.

\subsubsection{Почему нельзя писать \texttt{new T[newcap]}}

При реаллокации наивная реализация могла бы выглядеть так:

\begin{lstlisting}
T* newarr = new T[newcap]; // проблема!
\end{lstlisting}

Это не работает по двум причинам:
\begin{itemize}
    \item Мы не знаем, есть ли у \texttt{T} конструктор по умолчанию.
    \item Нам нужно выделить \textit{сырую} память и заполнять её по мере
          надобности, а не создавать все объекты сразу.
\end{itemize}

\subsubsection{Placement new}

Решение --- выделить сырую память через \texttt{char}, а объекты
создавать по месту с помощью \textbf{placement new}:

\begin{lstlisting}
new (ptr) T(args...);  // создаёт объект T по адресу ptr
\end{lstlisting}

Почему нельзя просто написать \texttt{newarr[i] = arr\_[i]}?
Потому что \texttt{newarr} указывает на сырые байты --- объекта
\texttt{T} там ещё нет, и присваивать можно только существующему объекту.
\texttt{memcpy} тоже не подходит: он не работает корректно, если объект
содержит указатель или ссылку на своё собственное поле
(например, \texttt{std::string} с small-string optimization или
декартово дерево).

\subsubsection{Почему нельзя \texttt{delete[] arr\_}}

При освобождении старой памяти нельзя писать \texttt{delete[] arr\_}:
это вызовет деструктор для каждого из \texttt{cap\_} слотов, хотя
реальных объектов там только \texttt{sz\_}. Нужно:
\begin{enumerate}
    \item вручную вызвать деструктор для каждого существующего объекта;
    \item освободить сырую память через
          \texttt{delete[] reinterpret\_cast<char*>(arr\_)}.
\end{enumerate}

\subsubsection{Реализация \texttt{reserve}}

\begin{lstlisting}
template<typename T>
class vector {
    T*     arr_;
    size_t sz_;
    size_t cap_;

    void reserve(size_t newcap) {
        T* newarr = reinterpret_cast<T*>(new char[newcap * sizeof(T)]);
        size_t i = 0;
        try {
            for (; i < sz_; ++i)
                new (newarr + i) T(arr_[i]); // placement new
        } catch (...) {
            for (size_t k = 0; k < i; ++k)
                (newarr + k)->~T();
            delete[] reinterpret_cast<char*>(newarr);
            throw;
        }
        for (size_t k = 0; k < sz_; ++k)
            (arr_ + k)->~T();
        delete[] reinterpret_cast<char*>(arr_);
        arr_ = newarr;
        cap_ = newcap;
    }
\end{lstlisting}

\textbf{Обработка исключений в \texttt{reserve}:}
\begin{itemize}
    \item Если \texttt{new char[...]} бросает \texttt{std::bad\_alloc}
          --- мы просто не выделили память, ничего не испорчено, вылетаем наверх.
    \item Если placement new бросает исключение в конструкторе на шаге \texttt{i}
          --- нужно уничтожить уже созданные объекты (с 0 по \texttt{i-1}),
          освободить \texttt{newarr} и перебросить исключение.
\end{itemize}

\subsubsection{Реализация \texttt{push\_back}}

\begin{lstlisting}
    void push_back(const T& value) {
        if (sz_ == cap_)
            reserve(cap_ > 0 ? cap_ * 2 : 1);
        new (arr_ + sz_) T(value);
        ++sz_;
    }
}; // конец class vector
\end{lstlisting}

% -------------------------------------------------------
\subsection{\texttt{vector<bool>}: специализация с битовым хранением}

\texttt{vector<bool>} --- это печально известная специализация,
которая ломает многие интуитивные ожидания от вектора.

\textbf{Идея:} хранить элементы побитово, упаковывая по 8 штук в \texttt{char}.
Это экономит память, но делает невозможным взять ссылку \texttt{bool\&} на
отдельный элемент --- ссылка не может указывать на отдельный бит.

\subsubsection{Прокси-класс \texttt{BitReference}}

Чтобы сохранить синтаксис \texttt{v[i] = true}, вводится прокси-класс:

\begin{lstlisting}
template<>
class vector<bool> {
    char*  arr_;
    size_t sz_;
    size_t cap_;

    struct BitReference {
        char*   cell;
        uint8_t idx;

        BitReference(char* cell, uint8_t idx) : cell(cell), idx(idx) {}

        void operator=(bool b) {
            if (b) *cell |=  (1 << idx);
            else   *cell &= ~(1 << idx);
        }

        operator bool() const {
            return *cell & (1 << idx);
        }
    };

public:
    BitReference operator[](size_t idx) {
        return BitReference(arr_ + idx / 8, idx % 8);
    }
};
\end{lstlisting}

\subsubsection{Трюк Ильи Мещерина: узнать тип выражения}

\begin{lstlisting}
template<typename T>
class Debug {
    Debug(T) = delete;
};

int main() {
    std::vector<bool> v(10);
    v[5] = true;
    // Debug d(v[5]); -- компилятор скажет тип: std::__bit_reference
}
\end{lstlisting}

\textbf{Почему это проблема:} \texttt{vector<bool>} нарушает концепцию
контейнера --- из него нельзя взять настоящую ссылку на элемент
(\texttt{bool\& r = v[5]} не скомпилируется), и многие шаблонные
алгоритмы ломаются при работе с ним. Поэтому для хранения битовых флагов
часто предпочитают \texttt{std::vector<char>} или \texttt{std::bitset}.

% -------------------------------------------------------
\newpage
\section{Библиотека \texttt{<algorithm>}}

\subsection{Зачем она вообще нужна}

Стандартная библиотека предоставляет \textbf{готовые алгоритмы},
которые работают с любыми контейнерами через \textbf{итераторы}.
Ключевая идея: алгоритм не знает, что за контейнер перед ним ---
он получает пару итераторов \texttt{[begin, end)} и работает с ними.

\begin{lstlisting}
#include <algorithm>
#include <vector>

int main() {
    int arr[] = {5, 3, 1, 4, 2};
    std::sort(arr, arr + 5);

    std::vector<int> v = {5, 3, 1, 4, 2};
    std::sort(v.begin(), v.end());
}
\end{lstlisting}

% -------------------------------------------------------
\subsection{\texttt{std::sort}}

Сортирует диапазон \textbf{на месте}. Сложность: $O(n \log n)$ (introsort).

\begin{lstlisting}
std::vector<int> v = {5, 3, 1, 4, 2};
std::sort(v.begin(), v.end()); // по возрастанию

std::sort(v.begin(), v.end(), [](int a, int b) {
    return a > b; // по убыванию
});
\end{lstlisting}

\textbf{Важно:} компаратор должен задавать \textit{строгий слабый порядок}
(\texttt{<}, а не \texttt{<=}). Если написать \texttt{a >= b} --- UB.

% -------------------------------------------------------
\subsection{\texttt{std::find} и \texttt{std::find\_if}}

Линейный поиск. Возвращает итератор на первый подходящий элемент,
или \texttt{last} если ничего не нашли.

\begin{lstlisting}
std::vector<int> v = {1, 2, 3, 4, 5};

auto it = std::find(v.begin(), v.end(), 3);
if (it != v.end()) std::cout << *it; // 3

auto it2 = std::find_if(v.begin(), v.end(), [](int x) {
    return x % 2 == 0; // первый чётный -> 2
});
\end{lstlisting}

\texttt{end()} как «не найдено» --- универсальное соглашение, потому что
алгоритм не знает о типе контейнера и не может вернуть \texttt{-1}.

% -------------------------------------------------------
\subsection{\texttt{std::count} и \texttt{std::count\_if}}

\begin{lstlisting}
std::vector<int> v = {1, 2, 2, 3, 2, 4};
int cnt  = std::count(v.begin(), v.end(), 2);                      // 3
int even = std::count_if(v.begin(), v.end(), [](int x){ return x%2==0; }); // 4
\end{lstlisting}

% -------------------------------------------------------
\subsection{\texttt{std::for\_each} и \texttt{std::transform}}

\begin{lstlisting}
std::vector<int> v = {1, 2, 3, 4, 5};

std::for_each(v.begin(), v.end(), [](int& x) { x *= 2; });
// v == {2, 4, 6, 8, 10}

std::vector<int> out(v.size());
std::transform(v.begin(), v.end(), out.begin(), [](int x) {
    return x * x;
});
// out == {4, 16, 36, 64, 100}
\end{lstlisting}

% -------------------------------------------------------
\subsection{\texttt{std::copy} и \texttt{std::back\_inserter}}

\begin{lstlisting}
std::vector<int> src = {1, 2, 3, 4, 5};
std::vector<int> evens;

std::copy_if(src.begin(), src.end(), std::back_inserter(evens),
             [](int x) { return x % 2 == 0; });
// evens == {2, 4}
\end{lstlisting}

\texttt{std::back\_inserter} --- итератор-адаптер, который при присваивании
вызывает \texttt{push\_back}. Нужен, когда выходной контейнер изначально пуст.

% -------------------------------------------------------
\subsection{Идиома erase-remove}

\texttt{std::remove} \textbf{не удаляет} элементы --- он лишь сдвигает
нужные в начало и возвращает итератор на новый логический конец.
Реальное удаление делает \texttt{erase}:

\begin{lstlisting}
std::vector<int> v = {1, 2, 3, 2, 4, 2, 5};
v.erase(std::remove(v.begin(), v.end(), 2), v.end());
// v == {1, 3, 4, 5}
\end{lstlisting}

% -------------------------------------------------------
\subsection{\texttt{std::accumulate} из \texttt{<numeric>}}

\begin{lstlisting}
#include <numeric>
std::vector<int> v = {1, 2, 3, 4, 5};
int sum     = std::accumulate(v.begin(), v.end(), 0);             // 15
int product = std::accumulate(v.begin(), v.end(), 1,
              [](int acc, int x) { return acc * x; });            // 120
\end{lstlisting}

% -------------------------------------------------------
\newpage
\section{Дек (\texttt{std::deque})}

\subsection{Инвалидация указателей и ссылок в \texttt{vector}}

\begin{lstlisting}
std::vector<int> v(10);
int* p = &v[5];
int& x = v[5];
v.push_back(1);      // реаллокация!
std::cout << *p;     // UB -- pointer invalidation
std::cout << x;      // UB -- reference invalidation
\end{lstlisting}

При \texttt{push\_back} вектор мог реаллоцироваться. Старый кусок
памяти нам больше не принадлежит --- все указатели и ссылки на элементы,
полученные до реаллокации, недействительны.

\subsection{Два главных отличия дека от вектора}

\begin{enumerate}
    \item \texttt{std::deque} \textbf{не инвалидирует} указатели и ссылки
          на существующие элементы при добавлении новых.
    \item Дек умеет \texttt{push\_front()} и \texttt{pop\_front()} за $O(1)$.
\end{enumerate}

\subsection{Внутреннее устройство дека}

Дек хранит \textbf{массив указателей на бакеты} фиксированного размера ---
\texttt{T** arr\_}.

\begin{itemize}
    \item При \texttt{push\_back}: идём вправо по бакету; когда бакет
          заполнен --- выделяем новый бакет справа.
    \item При \texttt{push\_front}: идём влево; когда бакет заполнен ---
          выделяем новый бакет слева и кладём элемент с его конца.
    \item Дек хранит две пары: (номер бакета начала, индекс внутри него)
          и (номер бакета конца, индекс внутри него).
\end{itemize}

\textbf{Почему не инвалидируются указатели?} При переполнении внешнего
массива реаллоцируется только он --- перепривязываются указатели на бакеты.
Сами бакеты с элементами остаются на месте.

\textbf{\texttt{operator[]} за $O(1)$:} зная индекс начала и размер бакета,
пересчитываем номер бакета и смещение за константное время.

\medskip
У дека нет \texttt{reserve()} и \texttt{shrink\_to\_fit()} --- элементы
разбросаны по бакетам, и их нельзя реаллоцировать одним куском.

% -------------------------------------------------------
\section{Адаптеры контейнеров}

\texttt{std::stack}, \texttt{std::queue} и \texttt{std::priority\_queue}
--- это \textbf{не контейнеры}, а \textbf{адаптеры}: они оборачивают
контейнер и предоставляют ограниченный интерфейс поверх него.

\begin{lstlisting}
std::stack<int>                    s;   // по умолчанию на deque
std::stack<int, std::vector<int>>  s2;  // на vector
std::stack<int, std::list<int>>    s3;  // на list
\end{lstlisting}

\subsection{\texttt{std::stack}}

\begin{lstlisting}
template<typename T, typename Container = std::deque<T>>
class stack {
    Container c_;
public:
    void   push(const T& val) { c_.push_back(val); }
    void   pop()              { c_.pop_back(); }
    T&     top()              { return c_.back(); }
    bool   empty() const      { return c_.empty(); }
    size_t size()  const      { return c_.size(); }
};
\end{lstlisting}

\textbf{Почему \texttt{pop()} возвращает \texttt{void}?}
Если бы \texttt{pop()} возвращал значение, пришлось бы скопировать
элемент перед удалением. Если конструктор копирования бросит исключение
--- элемент уже удалён и значение потеряно. Разделение на \texttt{top()}
(посмотреть) и \texttt{pop()} (удалить) даёт exception safety.
Вернуть ссылку тоже нельзя: после \texttt{pop()} она стала бы битой.

% -------------------------------------------------------
\newpage
\section{Итераторы и их категории}

\subsection{Что такое итератор}

Итератор --- это тип, который позволяет \textbf{обходить последовательность}.
Это обобщение указателя: его можно разыменовать (\texttt{*it}),
инкрементировать (\texttt{++it}) и сравнивать (\texttt{it != end}).
Обычный сырой указатель --- полноценный итератор.

Range-based for --- синтаксический сахар:
\begin{lstlisting}
for (int x : s) { ... }
// разворачивается в:
for (auto it = s.begin(); it != s.end(); ++it) { ... }
\end{lstlisting}
Достаточно реализовать \texttt{begin()} и \texttt{end()}.

\subsection{Категории итераторов}

\begin{center}
\begin{tabular}{lp{5.2cm}p{4cm}}
\toprule
\textbf{Категория} & \textbf{Что умеет} & \textbf{Примеры} \\
\midrule
\texttt{InputIterator}
    & \texttt{*it}, \texttt{++it}, \texttt{==}, \texttt{!=}
    & потоковые итераторы \\
\addlinespace
\texttt{ForwardIterator}
    & + гарантия повторного прохода
    & \texttt{unordered\_set}, \texttt{forward\_list} \\
\addlinespace
\texttt{BidirectionalIterator}
    & + \texttt{-{}-it}
    & \texttt{list}, \texttt{set}, \texttt{map} \\
\addlinespace
\texttt{RandomAccessIterator}
    & + \texttt{it+=n}, \texttt{it[n]}, \texttt{<}, \texttt{>}
    & \texttt{deque} \\
\addlinespace
\texttt{ContiguousIterator}
    & + гарантия непрерывности памяти
    & \texttt{vector}, \texttt{array}, указатель \\
\bottomrule
\end{tabular}
\end{center}

\medskip
Разница между \texttt{Input} и \texttt{Forward}: \texttt{InputIterator}
не гарантирует, что при повторном проходе увидим те же данные (поток ---
прочитанное не вернёшь). \texttt{ForwardIterator} эту гарантию даёт.

\subsection{\texttt{iterator\_traits}}

Для получения информации об итераторе внутри шаблонной функции:

\begin{lstlisting}
template<typename It>
void foo(It begin, It end) {
    typename std::iterator_traits<It>::value_type x = *begin;
    using Cat = typename std::iterator_traits<It>::iterator_category;
}
\end{lstlisting}

Доступные поля: \texttt{value\_type}, \texttt{reference}, \texttt{pointer},
\texttt{difference\_type}, \texttt{iterator\_category}.

Почему не \texttt{auto x = *begin}? Для \texttt{vector<bool>} \texttt{auto}
выведет \texttt{BitReference}, а не \texttt{bool}. \texttt{value\_type}
вернёт правильный тип.

\subsection{Реализация \texttt{std::distance}}

\begin{lstlisting}
template<typename It>
typename std::iterator_traits<It>::difference_type
distance(It first, It last) {
    if constexpr (std::is_same_v<
        typename std::iterator_traits<It>::iterator_category,
        std::random_access_iterator_tag>)
    {
        return last - first;  // O(1)
    }
    typename std::iterator_traits<It>::difference_type i = 0;
    for (; first != last; ++first, ++i) {}
    return i;                 // O(n)
}
\end{lstlisting}

\texttt{if constexpr} вычисляется на этапе компиляции: невыбранная ветка
не компилируется вообще. Это позволяет писать \texttt{last - first} не
боясь, что для non-random-access итератора это не скомпилируется.

\subsection{Реализация итератора вектора}

\begin{lstlisting}
template<typename T>
class Vector {
public:
    class Iterator {
    public:
        using iterator_category = std::random_access_iterator_tag;
        using value_type        = T;
        using difference_type   = std::ptrdiff_t;
        using pointer           = T*;
        using reference         = T&;

        explicit Iterator(pointer ptr) : ptr_(ptr) {}

        reference operator*()  const { return *ptr_; }
        pointer   operator->() const { return  ptr_; }

        Iterator& operator++()    { ++ptr_; return *this; }
        Iterator  operator++(int) { Iterator tmp = *this; ++ptr_; return tmp; }
        Iterator& operator--()    { --ptr_; return *this; }
        Iterator  operator--(int) { Iterator tmp = *this; --ptr_; return tmp; }

        Iterator  operator+(difference_type n) const { return Iterator(ptr_ + n); }
        Iterator  operator-(difference_type n) const { return Iterator(ptr_ - n); }
        difference_type operator-(const Iterator& o) const { return ptr_ - o.ptr_; }
        Iterator& operator+=(difference_type n) { ptr_ += n; return *this; }
        Iterator& operator-=(difference_type n) { ptr_ -= n; return *this; }
        reference operator[](difference_type n) const { return ptr_[n]; }

        bool operator==(const Iterator& o) const { return ptr_ == o.ptr_; }
        bool operator!=(const Iterator& o) const { return ptr_ != o.ptr_; }
        bool operator< (const Iterator& o) const { return ptr_ <  o.ptr_; }
        bool operator> (const Iterator& o) const { return ptr_ >  o.ptr_; }
        bool operator<=(const Iterator& o) const { return ptr_ <= o.ptr_; }
        bool operator>=(const Iterator& o) const { return ptr_ >= o.ptr_; }

    private:
        pointer ptr_;
    };

    Vector() : data_(nullptr), size_(0), capacity_(0) {}
    ~Vector() { delete[] data_; }

    void push_back(const T& val) {
        if (size_ == capacity_) reserve(capacity_ == 0 ? 1 : capacity_ * 2);
        data_[size_++] = val;
    }

    T&          operator[](std::size_t i) { return data_[i]; }
    std::size_t size()     const { return size_; }
    std::size_t capacity() const { return capacity_; }

    Iterator begin() { return Iterator(data_); }
    Iterator end()   { return Iterator(data_ + size_); }

private:
    void reserve(std::size_t new_cap) {
        T* tmp = new T[new_cap];
        for (std::size_t i = 0; i < size_; ++i) tmp[i] = data_[i];
        delete[] data_;
        data_ = tmp; capacity_ = new_cap;
    }

    T* data_; std::size_t size_; std::size_t capacity_;
};
\end{lstlisting}

Пять \texttt{type alias}-ов в \texttt{Iterator} --- контракт с
\texttt{iterator\_traits}. Без них \texttt{std::sort} и
\texttt{std::distance} не смогут узнать категорию итератора.

% -------------------------------------------------------
\section{Потоковые итераторы}

\subsection{\texttt{std::istream\_iterator}}

\begin{lstlisting}
#include <iterator>
#include <vector>

int main() {
    std::vector<int> v(
        std::istream_iterator<int>(std::cin),
        std::istream_iterator<int>()   // итератор-конец (EOF)
    );
}
\end{lstlisting}

Это \textbf{InputIterator} --- одноразовый: прочитанное назад не вернуть.
\texttt{istream\_iterator<int>()} без аргументов --- итератор конца потока.

\begin{lstlisting}
std::istream_iterator<int> in(std::cin), eof;
int sum = std::accumulate(in, eof, 0); // сумма всех чисел из stdin
\end{lstlisting}

\subsection{\texttt{std::ostream\_iterator}}

\begin{lstlisting}
std::vector<int> v = {1, 2, 3, 4, 5};

std::copy(v.begin(), v.end(),
          std::ostream_iterator<int>(std::cout, " "));
// вывод: 1 2 3 4 5

std::transform(v.begin(), v.end(),
               std::ostream_iterator<int>(std::cout, "\n"),
               [](int x) { return x * x; });
\end{lstlisting}

\texttt{ostream\_iterator} --- \textbf{OutputIterator}: присваивание
\texttt{*it = val} вызывает \texttt{os << val << delim}.

\subsection{Итераторы файловых потоков}

\begin{lstlisting}
#include <fstream>
#include <iterator>
#include <algorithm>

int main() {
    std::ifstream fin("input.txt");
    std::ofstream fout("output.txt");

    std::vector<int> v(std::istream_iterator<int>(fin),
                       std::istream_iterator<int>());
    std::sort(v.begin(), v.end());
    std::copy(v.begin(), v.end(),
              std::ostream_iterator<int>(fout, "\n"));
}
\end{lstlisting}

Алгоритмы не знают, откуда данные читаются и куда пишутся --- они
работают с абстракцией итератора. В этом и есть смысл всей иерархии.

% -------------------------------------------------------
\newpage
\section{Аллокаторы и управление памятью}

\subsection{Зачем аллокатор нужен контейнеру}

У каждого контейнера есть шаблонный параметр аллокатора:

\begin{lstlisting}
std::list<int>    // = std::list<int, std::allocator<int>>
std::vector<int>  // = std::vector<int, std::allocator<int>>
\end{lstlisting}

Аллокатор --- это \textbf{промежуточный класс между контейнером и
оператором \texttt{new}}. Стек вызовов:

\begin{center}
\texttt{Container} $\to$ \texttt{Allocator} $\to$
\texttt{operator new} $\to$ \texttt{malloc} $\to$ \texttt{OS}
\end{center}

\subsection{Устройство стандартного аллокатора}

У аллокатора четыре метода:

\begin{lstlisting}
template<typename T>
struct allocator {
    using value_type = T;

    T* allocate(std::size_t n) {
        if (n == 0) return nullptr;
        return static_cast<T*>(::operator new(n * sizeof(T)));
    }

    void deallocate(T* p, std::size_t) {
        ::operator delete(p);
    }

    template<typename... Args>
    void construct(T* p, Args&&... args) {
        ::new (static_cast<void*>(p)) T(std::forward<Args>(args)...);
    }

    void destroy(T* p) {
        p->~T();
    }
};
\end{lstlisting}

\begin{itemize}
    \item \texttt{allocate} --- выделяет сырую память, конструкторы не вызываются.
    \item \texttt{deallocate} --- освобождает память, деструкторы не вызываются.
    \item \texttt{construct} --- вызывает конструктор по адресу (placement new).
    \item \texttt{destroy} --- явно вызывает деструктор, память не освобождает.
\end{itemize}

\subsection{Проблема: \texttt{list} аллоцирует не \texttt{T}}

\texttt{std::list<int>} хранит не просто \texttt{int}, а узлы:

\begin{lstlisting}
template<typename T>
struct ListNode { T value; ListNode* prev; ListNode* next; };
\end{lstlisting}

Но аллокатор шаблонизирован по \texttt{T = int}. Решение комитета ---
\textbf{rebind}: вложенный шаблон в аллокаторе, который перепривязывает
его к другому типу:

\begin{lstlisting}
template<typename T>
struct allocator {
    template<typename U>
    struct rebind { using other = allocator<U>; };
};

// list внутри делает:
using NodeAlloc = typename Alloc::template rebind<ListNode<T>>::other;
NodeAlloc node_alloc_;  // аллокатор для узлов, а не для T
\end{lstlisting}

\subsection{\texttt{allocator\_traits}}

\texttt{std::allocator\_traits<Alloc>} --- адаптер с единым интерфейсом
к любому аллокатору и разумными дефолтами для отсутствующих методов.

\begin{lstlisting}
template<typename Alloc>
struct allocator_traits {
    using value_type      = typename Alloc::value_type;
    using pointer         = value_type*;
    using size_type       = std::size_t;
    using difference_type = std::ptrdiff_t;

    template<typename U>
    using rebind_alloc = typename Alloc::template rebind<U>::other;

    static pointer allocate(Alloc& a, size_type n)
        { return a.allocate(n); }
    static void deallocate(Alloc& a, pointer p, size_type n)
        { a.deallocate(p, n); }

    // если у Alloc нет construct -- placement new по умолчанию
    template<typename T, typename... Args>
    static void construct(Alloc&, T* p, Args&&... args)
        { ::new (static_cast<void*>(p)) T(std::forward<Args>(args)...); }

    // если у Alloc нет destroy -- явный деструктор по умолчанию
    template<typename T>
    static void destroy(Alloc&, T* p) { p->~T(); }

    // политики распространения аллокатора
    using propagate_on_container_copy_assignment = std::false_type;
    using propagate_on_container_move_assignment = std::false_type;
    using propagate_on_container_swap            = std::false_type;
};
\end{lstlisting}

\textbf{Главная идея}: контейнер должен работать с аллокатором
\textit{только через \texttt{allocator\_traits}}. Тогда кастомный
аллокатор может не реализовывать \texttt{construct}/\texttt{destroy}
и всё равно работать.

\subsubsection{Политики распространения аллокатора}

\begin{center}
\begin{tabular}{lp{7.5cm}}
\toprule
\textbf{Флаг} & \textbf{Что означает} \\
\midrule
\texttt{propagate\_on\_container\_copy\_assignment}
    & Копировать аллокатор при \texttt{operator=} копирования? \\
\addlinespace
\texttt{propagate\_on\_container\_move\_assignment}
    & Копировать аллокатор при \texttt{operator=} перемещения? \\
\addlinespace
\texttt{propagate\_on\_container\_swap}
    & Менять аллокаторы местами при \texttt{swap}? \\
\bottomrule
\end{tabular}
\end{center}

По умолчанию все три --- \texttt{false\_type}. Это важно для stateful
аллокаторов: если два вектора используют разные арены, при свопе нужно
либо перетащить элементы, либо поменять аллокаторы.

\subsection{Allocator-aware контейнер}

\begin{lstlisting}
template<typename T, typename Alloc = std::allocator<T>>
class vector {
    using AllocTraits = std::allocator_traits<Alloc>;

    Alloc  alloc_;
    T*     arr_;
    size_t sz_;
    size_t cap_;

public:
    explicit vector(const Alloc& a = Alloc())
        : alloc_(a), arr_(nullptr), sz_(0), cap_(0) {}

    void push_back(const T& val) {
        if (sz_ == cap_) reserve(cap_ ? cap_ * 2 : 1);
        AllocTraits::construct(alloc_, arr_ + sz_, val);
        ++sz_;
    }

    ~vector() {
        for (size_t i = 0; i < sz_; ++i)
            AllocTraits::destroy(alloc_, arr_ + i);
        AllocTraits::deallocate(alloc_, arr_, cap_);
    }

private:
    void reserve(size_t newcap) {
        T* newarr = AllocTraits::allocate(alloc_, newcap);
        size_t i = 0;
        try {
            for (; i < sz_; ++i)
                AllocTraits::construct(alloc_, newarr + i, arr_[i]);
        } catch (...) {
            for (size_t k = 0; k < i; ++k)
                AllocTraits::destroy(alloc_, newarr + k);
            AllocTraits::deallocate(alloc_, newarr, newcap);
            throw;
        }
        for (size_t k = 0; k < sz_; ++k)
            AllocTraits::destroy(alloc_, arr_ + k);
        AllocTraits::deallocate(alloc_, arr_, cap_);
        arr_ = newarr; cap_ = newcap;
    }
};
\end{lstlisting}

Stateless аллокатор (как \texttt{std::allocator}) не занимает места
благодаря \textbf{EBO} (Empty Base Optimization).

\subsection{Перегрузка \texttt{operator new} и \texttt{operator delete}}

\subsubsection{Глобальная перегрузка}

\begin{lstlisting}
void* operator new(std::size_t size) {
    void* ptr = std::malloc(size);
    if (!ptr) throw std::bad_alloc{};
    return ptr;
}
void operator delete(void* ptr) noexcept { std::free(ptr); }

// C++14: sized delete
void operator delete(void* ptr, std::size_t size) noexcept { std::free(ptr); }
\end{lstlisting}

Затрагивает весь процесс. Используется для профилировщиков памяти и
детекторов утечек.

\subsubsection{Классовая перегрузка}

\begin{lstlisting}
struct Node {
    int value; Node* next;

    void* operator new(std::size_t size) {
        return ::operator new(size); // делегируем глобальному
    }
    void operator delete(void* ptr) noexcept {
        ::operator delete(ptr);
    }
};
\end{lstlisting}

При \texttt{new Node} компилятор сначала ищет \texttt{operator new}
в классе, и только если не нашёл --- берёт глобальный.

\subsubsection{Когда использовать что}

\begin{center}
\begin{tabular}{lp{7.5cm}}
\toprule
\textbf{Инструмент} & \textbf{Когда применять} \\
\midrule
Кастомный аллокатор
    & Своя стратегия памяти для конкретного контейнера \\
\addlinespace
Классовая перегрузка \texttt{new/delete}
    & Свой пул для конкретного класса (например, узлы дерева) \\
\addlinespace
Глобальная перегрузка \texttt{new/delete}
    & Профилирование, детектирование утечек для всей программы \\
\bottomrule
\end{tabular}
\end{center}

% -------------------------------------------------------
\newpage
\section{Многопоточность в C++}

\begin{center}
\fbox{\parbox{0.9\textwidth}{
\textbf{Лирическое отступление.} Данный раздел \textbf{не относится} к текущему
курсу и приводится исключительно в ознакомительных целях. Многопоточность будет
подробно рассмотрена в последующих курсах.
}}
\end{center}

\medskip

% -------------------------------------------------------
\subsection{Определение потока исполнения}

\subsubsection{Что такое поток?}

\textbf{Поток исполнения (thread of execution)} --- это единственная последовательная
цепочка команд, выполняемых процессором. Представьте это как <<линию чтения>> в книге:
один палец ведёт по одной строке. Если у вас два пальца --- вы читаете два куска книги
одновременно (два потока).

Программа считается завершённой \textbf{только тогда}, когда завершились \textbf{все}
её потоки. Даже если \texttt{main()} уже отработал, отсоединённые потоки могут ещё работать.

\subsubsection{Аналогия}

Представьте кухню ресторана:
\begin{itemize}
    \item \textbf{Однопоточная программа} --- один повар делает всё по очереди:
          режет овощи, варит суп, жарит мясо. Пока суп варится, повар просто ждёт.
    \item \textbf{Многопоточная программа} --- несколько поваров работают параллельно.
          Один варит суп, другой жарит мясо, третий режет салат.
\end{itemize}

\subsubsection{Логическая многопоточность vs Аппаратный параллелизм}

Это \textbf{не одно и то же}:

\begin{center}
\begin{tabular}{lp{5cm}p{5cm}}
\toprule
\textbf{Понятие} & \textbf{Суть} & \textbf{Пример} \\
\midrule
Логическая многопоточность
    & ОС быстро переключается между потоками на одном ядре, создавая \textit{иллюзию} параллельности
    & 1 ядро, 4 потока --- ОС по очереди даёт каждому \textasciitilde10 мс \\
\addlinespace
Аппаратный параллелизм
    & Потоки реально выполняются одновременно на разных ядрах
    & 4 ядра $\times$ 1 поток --- каждый работает физически одновременно \\
\bottomrule
\end{tabular}
\end{center}

\medskip
Многопоточность появилась \textbf{раньше} многоядерных процессоров. Зачем?
Чтобы решить проблему \textbf{блокирующего ввода-вывода}. Классический пример:
в ранних ОС форматирование дискеты <<вешало>> всю систему, потому что процессор
просто ждал завершения I/O-операции. С появлением потоков ОС могла переключиться
на другую задачу, пока диск форматируется.

\subsubsection{Типы многозадачности}

\textbf{Кооперативная многозадачность} --- поток \textit{сам} решает, когда
отдать управление другим. Если поток <<зажадничал>> и не отдаёт CPU ---
все остальные замирают.

\begin{lstlisting}[numbers=none]
Поток A: работаю... работаю... ладно, передаю управление -> yield()
Поток B: о, моя очередь! работаю... yield()
Поток A: снова я!
\end{lstlisting}

Пример в современном C++ --- \textbf{корутины} (\texttt{co\_await},
\texttt{co\_yield} в C++20).

\medskip
\textbf{Вытесняющая многозадачность} --- ОС \textit{сама} прерывает потоки
по таймеру (обычно каждые \textasciitilde10 мс). Поток не может <<зависнуть>>
и заблокировать всю систему.

\begin{lstlisting}[numbers=none]
Поток A: работаю... работ--
[ОС]: Время вышло! Переключаю на B.
Поток B: работаю...
\end{lstlisting}

Это стандарт для всех современных ОС (Windows, Linux, macOS).

% -------------------------------------------------------
\subsection{Работа с потоками в C++}

\subsubsection{Создание потока}

Потоки создаются через \texttt{std::thread} (доступен с C++11). Важная деталь:
\textbf{поток запускается немедленно} в момент создания объекта --- нет отдельного
метода \texttt{start()}.

\begin{lstlisting}
#include <iostream>
#include <thread>

void say_hello() {
    std::cout << "Привет из потока!\n";
}

int main() {
    std::thread t(say_hello);  // Поток УЖЕ работает!
    t.join();                  // Ждём завершения
    return 0;
}
\end{lstlisting}

\subsubsection{Передача аргументов в поток}

\begin{lstlisting}
void print_sum(int a, int b) {
    std::cout << a + b << "\n";
}

int main() {
    std::thread t(print_sum, 3, 7);  // Аргументы передаются после функции
    t.join();
}
\end{lstlisting}

\textbf{Внимание:} аргументы \textbf{копируются} в поток.
Если нужна ссылка, используйте \texttt{std::ref()}:

\begin{lstlisting}
void increment(int& x) { x++; }

int main() {
    int value = 0;
    std::thread t(increment, std::ref(value));  // Передаём ссылку
    t.join();
    std::cout << value;  // 1
}
\end{lstlisting}

\subsubsection{Использование лямбд}

\begin{lstlisting}
int main() {
    int data = 42;
    std::thread t([&data]() {
        data *= 2;
        std::cout << "В потоке: " << data << "\n";
    });
    t.join();
    std::cout << "В main: " << data << "\n";  // 84
}
\end{lstlisting}

\subsubsection{\texttt{join()} vs \texttt{detach()}}

\begin{center}
\begin{tabular}{lp{5cm}p{5cm}}
\toprule
\textbf{Метод} & \textbf{Что делает} & \textbf{Когда использовать} \\
\midrule
\texttt{join()}
    & Блокирует текущий поток до завершения вызванного
    & Когда нужен результат работы потока \\
\addlinespace
\texttt{detach()}
    & Отсоединяет поток, он живёт самостоятельно
    & Фоновые задачи (логирование, мониторинг) \\
\bottomrule
\end{tabular}
\end{center}

\begin{lstlisting}
// Пример join() -- ждём результат
std::thread t(heavy_computation);
// ... тут main делает свои дела ...
t.join();  // Теперь точно результат готов

// Пример detach() -- fire and forget
std::thread logger(write_logs);
logger.detach();  // Пусть пишет логи в фоне
\end{lstlisting}

\subsubsection{Опасность: забыть \texttt{join}/\texttt{detach}}

Если объект \texttt{std::thread} уничтожается без вызова \texttt{join()} или
\texttt{detach()}, программа \textbf{аварийно завершается} через
\texttt{std::terminate()}:

\begin{lstlisting}
void dangerous() {
    std::thread t(some_work);
    // Упс! Функция завершилась, t уничтожается
    // без join/detach -> std::terminate()!
}
\end{lstlisting}

\textbf{Безопасный паттерн} --- RAII-обёртка (или \texttt{std::jthread} из C++20,
который сам делает \texttt{join} в деструкторе):

\begin{lstlisting}
// C++20
void safe() {
    std::jthread t(some_work);
    // Деструктор jthread сам вызовет join()
}
\end{lstlisting}

% -------------------------------------------------------
\subsection{Область памяти и состояние гонки (Data Race)}

\subsubsection{Что такое Memory Location?}

\textbf{Memory Location} (область памяти) по стандарту --- это:
\begin{itemize}
    \item Объект \textbf{скалярного типа} (\texttt{int}, \texttt{char},
          \texttt{float}, указатель и~т.\,д.)
    \item Или \textbf{максимальная последовательность битовых полей ненулевой длины}
\end{itemize}

Важный нюанс: два соседних поля структуры --- это \textbf{разные} области памяти,
к ним можно обращаться параллельно без гонки. Но соседние битовые поля --- это
\textbf{одна} область памяти!

\begin{lstlisting}
struct S {
    int x;      // Отдельная memory location
    int y;      // Отдельная memory location -- можно менять параллельно с x
};

struct Bits {
    unsigned a : 4;  // a и b -- ОДНА memory location!
    unsigned b : 4;  // Параллельное изменение a и b = Data Race
    unsigned   : 0;  // Разделитель
    unsigned c : 4;  // c -- отдельная memory location
};
\end{lstlisting}

\subsubsection{Что такое Data Race?}

\textbf{Data Race (гонка данных)} --- это ситуация, когда:

\begin{enumerate}
    \item Два потока обращаются к \textbf{одной области памяти}
    \item \textbf{Хотя бы один} из них \textbf{записывает}
    \item Между ними \textbf{нет синхронизации} (нет happens-before)
    \item Хотя бы одно действие \textbf{не атомарное}
\end{enumerate}

\textbf{Data Race = Undefined Behavior (UB).} Это не просто <<неправильный результат>>.
Компилятор имеет право на \textit{что угодно}: удалить код, выдать мусор,
отформатировать диск (теоретически).

\subsubsection{Наглядный пример гонки}

\begin{lstlisting}
#include <thread>

int counter = 0;  // Общая переменная

void increment_1000() {
    for (int i = 0; i < 1000; ++i) {
        counter++;  // DATA RACE!
    }
}

int main() {
    std::thread t1(increment_1000);
    std::thread t2(increment_1000);
    t1.join();
    t2.join();

    // Ожидаем 2000, но результат непредсказуем!
    // Может быть 1500, 1837, 2000, или что угодно
    std::cout << counter << "\n";
}
\end{lstlisting}

\textbf{Почему не 2000?} Операция \texttt{counter++} не атомарна.
На уровне процессора это три шага:

\begin{lstlisting}[numbers=none]
1. Прочитать counter из памяти -> регистр    (LOAD)
2. Увеличить значение в регистре на 1       (ADD)
3. Записать результат обратно в память       (STORE)
\end{lstlisting}

Если два потока выполняют это одновременно:

\begin{lstlisting}[numbers=none]
Поток 1: LOAD counter (=0) -> регистр = 0
Поток 2: LOAD counter (=0) -> регистр = 0
Поток 1: ADD -> регистр = 1
Поток 2: ADD -> регистр = 1
Поток 1: STORE -> counter = 1
Поток 2: STORE -> counter = 1   <- Инкремент потерян!
\end{lstlisting}

% -------------------------------------------------------
\subsection{Распространённые ошибки и заблуждения}

\subsubsection{Ошибка №1: \texttt{volatile} <<чинит>> гонки}

\textbf{Нет!} \texttt{volatile} --- это указание компилятору:
\begin{itemize}
    \item Не оптимизировать чтение/запись этой переменной
    \item Каждый раз читать значение из памяти, а не из регистра
\end{itemize}

Но \texttt{volatile} \textbf{НЕ обеспечивает}:
\begin{itemize}
    \item Атомарность операций
    \item Барьеры памяти (memory barriers)
    \item Когерентность кэшей между ядрами
\end{itemize}

\begin{lstlisting}
volatile int counter = 0;  // НЕ решает проблему!

// counter++ по-прежнему = LOAD + ADD + STORE
// volatile лишь гарантирует, что компилятор не закэширует
// значение в регистре, но гонка всё равно есть

// Правильное решение:
std::atomic<int> counter = 0;  // Атомарная переменная
\end{lstlisting}

\texttt{volatile} полезен для:
\begin{itemize}
    \item Регистров аппаратуры (MMIO)
    \item Переменных, изменяемых обработчиками сигналов
    \item Переменных, изменяемых через \texttt{setjmp}/\texttt{longjmp}
\end{itemize}

\subsubsection{Ошибка №2: Привязка к одному ядру решает гонки}

Можно использовать системные вызовы для привязки потоков к одному ядру:

\begin{lstlisting}
// Linux-специфичный код
cpu_set_t cpuset;
CPU_ZERO(&cpuset);
CPU_SET(0, &cpuset);  // Привязать к ядру 0
sched_setaffinity(0, sizeof(cpuset), &cpuset);
\end{lstlisting}

Визуально симптомы гонки \textbf{могут исчезнуть} (потоки не работают
\textit{физически} параллельно). Но по стандарту C++ это \textbf{всё равно UB},
потому что:
\begin{itemize}
    \item Язык C++ <<не знает>> про ядра процессора
    \item Стандарт определяет поведение абстрактной машины, а не конкретного железа
    \item Компилятор может оптимизировать код, опираясь на правила стандарта
\end{itemize}

\subsubsection{Ошибка №3: \texttt{std::cout} безопасен для многопоточности}

Стандарт \textbf{гарантирует}, что в \texttt{std::cout} не будет Data Race
(внутренний буфер защищён). Но \textbf{порядок вывода} может перемешиваться:

\begin{lstlisting}
// Поток 1:
std::cout << "Hello " << "from " << "thread 1\n";

// Поток 2:
std::cout << "Hello " << "from " << "thread 2\n";

// Возможный вывод:
// Hello Hello from from thread 1thread 2
//
// Каждая отдельная операция << атомарна,
// но между ними другой поток может вклиниться
\end{lstlisting}

\textbf{Решение --- вывод одной строкой:}

\begin{lstlisting}
// Вариант 1: формируем строку заранее
std::string msg = "Hello from thread 1\n";
std::cout << msg;  // Один вызов << = один непрерывный вывод

// Вариант 2: std::osyncstream (C++20)
std::osyncstream(std::cout) << "Hello " << "from " << "thread 1\n";
// Весь вывод буферизуется и отправляется атомарно
\end{lstlisting}

% -------------------------------------------------------
\subsection{Синхронизация: Мьютексы}

\subsubsection{Что такое мьютекс?}

\textbf{Мьютекс (mutex = mutual exclusion)} --- это объект синхронизации,
который гарантирует, что только \textbf{один поток} в каждый момент времени
выполняет защищённый участок кода (критическую секцию).

Аналогия: туалетная кабинка с замком. Зашёл --- закрыл (\texttt{lock}).
Вышел --- открыл (\texttt{unlock}). Если занято --- жди.

\subsubsection{Базовый интерфейс}

\begin{lstlisting}
#include <mutex>

std::mutex mtx;
int counter = 0;

void safe_increment() {
    for (int i = 0; i < 1000; ++i) {
        mtx.lock();       // Захватить мьютекс (или ждать)
        counter++;        // Критическая секция -- только один поток
        mtx.unlock();     // Освободить мьютекс
    }
}

int main() {
    std::thread t1(safe_increment);
    std::thread t2(safe_increment);
    t1.join();
    t2.join();
    std::cout << counter << "\n";  // Всегда 2000!
}
\end{lstlisting}

\subsubsection{Методы мьютекса}

\begin{center}
\begin{tabular}{lp{9cm}}
\toprule
\textbf{Метод} & \textbf{Поведение} \\
\midrule
\texttt{lock()}
    & Захватить. Если занят --- \textbf{заблокировать} текущий поток до освобождения \\
\addlinespace
\texttt{unlock()}
    & Освободить мьютекс \\
\addlinespace
\texttt{try\_lock()}
    & Попробовать захватить. Если занят --- \textbf{не ждать}, вернуть \texttt{false} \\
\bottomrule
\end{tabular}
\end{center}

\begin{lstlisting}
if (mtx.try_lock()) {
    // Захватили! Делаем работу...
    mtx.unlock();
} else {
    // Мьютекс занят, делаем что-то другое
}
\end{lstlisting}

\subsubsection{Правила и UB при работе с мьютексами}

\textbf{UB №1: Повторный захват в том же потоке (рекурсивный lock)}

\begin{lstlisting}
std::mutex mtx;

void foo() {
    mtx.lock();
    // ... делаем что-то ...
    mtx.lock();   // UB! Этот поток уже владеет мьютексом!
                  // Deadlock или неопределённое поведение
    mtx.unlock();
    mtx.unlock();
}
\end{lstlisting}

Решение --- \texttt{std::recursive\_mutex} (если рекурсия действительно нужна):

\begin{lstlisting}
std::recursive_mutex rmtx;

void foo() {
    rmtx.lock();
    rmtx.lock();   // Рекурсивный мьютекс допускает это
    rmtx.unlock();
    rmtx.unlock(); // Нужно столько же unlock, сколько lock
}
\end{lstlisting}

\textbf{UB №2: Разблокировка чужого мьютекса}

\begin{lstlisting}
std::mutex mtx;

// Поток 1:
mtx.lock();
// работает...

// Поток 2:
mtx.unlock();  // UB! Поток 2 не захватывал этот мьютекс!
\end{lstlisting}

\textbf{Почему нет метода \texttt{is\_locked()}?}

Потому что он бесполезен: к моменту, когда вы получите ответ, состояние уже
могло измениться:

\begin{lstlisting}
// Бессмысленный гипотетический код:
if (!mtx.is_locked()) {    // В этот момент мьютекс свободен
    mtx.lock();             // Но здесь другой поток мог его захватить!
}
// Именно поэтому существует try_lock() -- атомарная проверка+захват
\end{lstlisting}

% -------------------------------------------------------
\subsection{RAII в многопоточности}

\subsubsection{Проблема голых \texttt{lock}/\texttt{unlock}}

\begin{lstlisting}
void dangerous_function() {
    mtx.lock();

    do_work();          // Что если тут исключение?

    if (error) {
        return;         // Забыли unlock! Мьютекс навсегда заблокирован!
    }

    mtx.unlock();       // До сюда можем не дойти
}
\end{lstlisting}

\subsubsection{\texttt{std::lock\_guard} --- простая RAII-обёртка}

\textbf{Принцип RAII:} конструктор захватывает ресурс, деструктор освобождает.
Гарантирует освобождение при \textbf{любом} выходе из области видимости
(\texttt{return}, исключение, \texttt{break}\ldots).

\begin{lstlisting}
void safe_function() {
    std::lock_guard<std::mutex> lock(mtx);  // lock() в конструкторе

    do_work();          // Исключение? Не страшно!

    if (error) {
        return;         // Деструктор lock_guard вызовет unlock()
    }

    // Деструктор lock_guard вызовет unlock()
}
\end{lstlisting}

С C++17 можно не указывать тип (CTAD):

\begin{lstlisting}
std::lock_guard lock(mtx);  // Компилятор выведет тип сам
\end{lstlisting}

\subsubsection{\texttt{std::unique\_lock} --- гибкая обёртка}

В отличие от \texttt{lock\_guard}, позволяет:
\begin{itemize}
    \item Откладывать захват (\texttt{std::defer\_lock})
    \item Досрочно освобождать (\texttt{unlock()})
    \item Перемещать между областями видимости
\end{itemize}

\begin{lstlisting}
void flexible_function() {
    std::unique_lock<std::mutex> lock(mtx, std::defer_lock);
    // Мьютекс ещё НЕ захвачен

    prepare_data();

    lock.lock();          // Захватываем, когда нужно
    modify_shared_data();
    lock.unlock();        // Освобождаем досрочно

    do_unrelated_work();  // Работаем без мьютекса
}
\end{lstlisting}

\subsubsection{\texttt{std::scoped\_lock} (C++17) --- захват нескольких мьютексов}

Главное преимущество: захватывает \textbf{несколько} мьютексов \textbf{без риска дедлока}:

\begin{lstlisting}
std::mutex mtx1, mtx2;

void transfer(Account& from, Account& to, int amount) {
    std::scoped_lock lock(mtx1, mtx2);  // Оба мьютекса захвачены безопасно
    from.balance -= amount;
    to.balance += amount;
    // Деструктор освободит оба мьютекса
}
\end{lstlisting}

\subsubsection{Сравнение RAII-обёрток}

\begin{center}
\begin{tabular}{lcccc}
\toprule
\textbf{Свойство} & \texttt{lock\_guard} & \texttt{unique\_lock} & \texttt{scoped\_lock} \\
\midrule
Несколько мьютексов   & нет & нет & да  \\
Отложенный захват     & нет & да  & нет \\
Досрочное освобождение & нет & да  & нет \\
Перемещение           & нет & да  & нет \\
Стандарт              & C++11 & C++11 & C++17 \\
\bottomrule
\end{tabular}
\end{center}

% -------------------------------------------------------
\subsection{Проектирование потокобезопасных объектов}

\subsubsection{Уровни потокобезопасности}

\textbf{Уровень <<Нейтральный>>} --- объект сам по себе не защищён, но безопасен
при внешней синхронизации.

\begin{lstlisting}
int x = 0;  // Нейтральный тип

// Безопасно, если обернуть в мьютекс:
std::mutex mtx;
{
    std::lock_guard lock(mtx);
    x++;
}
\end{lstlisting}

\textbf{Уровень <<Ниже нейтрального>>} --- указатели. Даже если синхронизировать
сам указатель, данные по нему могут быть не защищены:

\begin{lstlisting}
std::mutex mtx;
int* ptr = new int(42);

// Поток 1:
{
    std::lock_guard lock(mtx);
    ptr = new int(100);       // Защитили ptr
}

// Поток 2:
int local = *ptr;             // А данные по указателю -- не защищены!
                              // ptr мог уже указывать на удалённую память
\end{lstlisting}

\textbf{Уровень <<Потокобезопасный>>} --- все операции внутренне защищены.
Пример: \texttt{std::atomic<int>}.

\begin{lstlisting}
std::atomic<int> counter = 0;

// Поток 1:
counter++;  // Атомарно и безопасно

// Поток 2:
counter++;  // Атомарно и безопасно
\end{lstlisting}

\subsubsection{API Race --- коварная ловушка}

Это ситуация, когда \textbf{каждый отдельный метод} защищён мьютексом
(Data Race нет!), но \textbf{последовательность вызовов} создаёт логическую ошибку.

\textbf{Классический пример: потокобезопасный стек}

\begin{lstlisting}
// Допустим, у нас есть "потокобезопасный" стек:
class ThreadSafeStack {
    std::mutex mtx;
    std::stack<int> data;
public:
    bool empty() {
        std::lock_guard lock(mtx);
        return data.empty();
    }
    int top() {
        std::lock_guard lock(mtx);
        return data.top();
    }
    void pop() {
        std::lock_guard lock(mtx);
        data.pop();
    }
};
\end{lstlisting}

Каждый метод по отдельности безопасен. Но посмотрите:

\begin{lstlisting}
ThreadSafeStack stack;

// Поток 1:
if (!stack.empty()) {      // 1. Проверили -- не пуст
                           // Между empty() и top() поток 2 вызвал pop()!
    int val = stack.top(); // 2. Стек уже пуст -> UB!
    stack.pop();
}

// Поток 2 (выполняется между шагами 1 и 2 потока 1):
if (!stack.empty()) {
    stack.pop();           // Забрал последний элемент
}
\end{lstlisting}

\textbf{Решение: объединение операций}

\begin{lstlisting}
class SafeStack {
    std::mutex mtx;
    std::stack<int> data;
public:
    // Одна атомарная операция вместо трёх
    std::optional<int> try_pop() {
        std::lock_guard lock(mtx);
        if (data.empty()) {
            return std::nullopt;  // Пусто -- вернули "ничего"
        }
        int val = data.top();
        data.pop();
        return val;               // Проверка + чтение + удаление = одна операция
    }

    void push(int val) {
        std::lock_guard lock(mtx);
        data.push(val);
    }
};

// Использование:
if (auto val = stack.try_pop()) {
    std::cout << "Получили: " << *val << "\n";
} else {
    std::cout << "Стек пуст\n";
}
\end{lstlisting}

% -------------------------------------------------------
\subsection{Дедлоки (Deadlocks)}

\subsubsection{Что такое дедлок?}

\textbf{Дедлок (взаимная блокировка)} --- ситуация, когда два (или более) потока
бесконечно ждут друг друга, потому что каждый держит ресурс, нужный другому.

\subsubsection{Классический пример: перевод денег}

\begin{lstlisting}
std::mutex mtx_A, mtx_B;  // Мьютексы для счетов A и B

// Поток 1: перевод A -> B
void transfer_A_to_B() {
    mtx_A.lock();           // 1. Захватил A
    // ... тут переключение контекста ...
    mtx_B.lock();           // 3. Ждёт B... но B занят потоком 2 -> ДЕДЛОК
    // перевод...
    mtx_B.unlock();
    mtx_A.unlock();
}

// Поток 2: перевод B -> A
void transfer_B_to_A() {
    mtx_B.lock();           // 2. Захватил B
    // ... тут переключение контекста ...
    mtx_A.lock();           // 4. Ждёт A... но A занят потоком 1 -> ДЕДЛОК
    // перевод...
    mtx_A.unlock();
    mtx_B.unlock();
}
\end{lstlisting}

Визуально:

\begin{lstlisting}[numbers=none]
Поток 1: держит [A], хочет [B] ------+
                                      +-- Оба ждут вечно!
Поток 2: держит [B], хочет [A] ------+
\end{lstlisting}

\subsubsection{Задача об обедающих философах}

5 философов сидят за круглым столом. Между каждыми двумя лежит одна вилка.
Чтобы есть, нужны две вилки. Если каждый одновременно возьмёт левую вилку ---
все будут ждать правую вечно = дедлок.

\subsubsection{Решения в C++}

\textbf{Решение 1: \texttt{std::lock()} --- алгоритм безопасного захвата}

\begin{lstlisting}
std::mutex mtx_A, mtx_B;

void safe_transfer() {
    std::lock(mtx_A, mtx_B);  // Захватывает оба без дедлока!
    // std::lock использует внутренний алгоритм:
    // try_lock на одном -> try_lock на другом -> если не получилось,
    // отпустить всё и попробовать заново

    // Оборачиваем в lock_guard с adopt_lock (мьютексы уже захвачены)
    std::lock_guard lockA(mtx_A, std::adopt_lock);
    std::lock_guard lockB(mtx_B, std::adopt_lock);

    // Делаем перевод...
    // Деструкторы lock_guard освободят мьютексы
}
\end{lstlisting}

\textbf{Решение 2: \texttt{std::scoped\_lock} (C++17) --- всё в одну строку}

\begin{lstlisting}
void safe_transfer_modern() {
    std::scoped_lock lock(mtx_A, mtx_B);  // Всё! Безопасно и RAII
    // Делаем перевод...
}
\end{lstlisting}

\subsubsection{4 необходимых условия дедлока (условия Коффмана)}

Дедлок возможен \textbf{только} если выполнены \textbf{все четыре} условия одновременно:

\begin{enumerate}
    \item \textbf{Взаимное исключение} --- ресурс может быть захвачен только одним потоком
    \item \textbf{Удержание и ожидание} --- поток держит один ресурс и ждёт другой
    \item \textbf{Невытесняемость} --- ресурс нельзя отобрать у потока силой
    \item \textbf{Циклическое ожидание} --- замкнутая цепочка потоков, каждый ждёт следующего
\end{enumerate}

Разрушение \textbf{любого} из условий предотвращает дедлок. \texttt{std::lock}
разрушает условие №2 (захватывает всё или ничего).

% -------------------------------------------------------
\subsection{Итого: шпаргалка}

\begin{center}
\begin{tabular}{lp{8cm}}
\toprule
\textbf{Проблема} & \textbf{Решение} \\
\midrule
Забыл \texttt{join}/\texttt{detach}
    & \texttt{std::jthread} (C++20) \\
\addlinespace
Data Race
    & \texttt{std::mutex} + RAII-обёртки или \texttt{std::atomic} \\
\addlinespace
Забыл \texttt{unlock}
    & \texttt{std::lock\_guard} / \texttt{std::scoped\_lock} \\
\addlinespace
API Race
    & Объединение операций (\texttt{try\_pop} и~т.\,п.) \\
\addlinespace
Дедлок
    & \texttt{std::scoped\_lock(mtx1, mtx2, ...)} \\
\addlinespace
\texttt{volatile} вместо \texttt{atomic}
    & Заменить на \texttt{std::atomic<T>} \\
\addlinespace
Перемешанный вывод
    & \texttt{std::osyncstream} (C++20) \\
\bottomrule
\end{tabular}
\end{center}

% -------------------------------------------------------
\newpage
\section{Приведение типов (касты) в C++}

% -------------------------------------------------------
\subsection{<<Запрещённые заклинания>> --- зачем нужны касты?}

В C++ система типов --- это ваш главный защитник. Она ловит ошибки на этапе
компиляции: нельзя случайно сложить строку с указателем или передать \texttt{float}
туда, где ожидается \texttt{std::vector}. Но иногда программисту \textbf{нужно}
сказать компилятору: <<Я знаю, что делаю --- преобразуй этот тип в другой>>.

Для этого существуют \textbf{четыре заклинания} (каста), каждое со своей силой
и опасностью. Лектор не зря называет их <<запрещёнными>> --- неправильное
использование ведёт к Undefined Behavior.

\subsubsection{Почему не один каст, а четыре?}

В языке C был всего один способ: \texttt{(тип)выражение}. Он делал \textit{всё
подряд} --- и безопасные, и опасные преобразования --- и программист не мог отличить
одно от другого. C++ разделил эту грубую силу на четыре точных инструмента, чтобы:

\begin{itemize}
    \item Код \textbf{явно показывал намерение} программиста
    \item Опасные операции \textbf{бросались в глаза} при чтении кода
    \item Можно было \textbf{искать} конкретные касты grep'ом (попробуйте найти
          \texttt{(int)} в миллионе строк --- удачи)
\end{itemize}

% -------------------------------------------------------
\subsection{C-style cast --- <<Грубая сила>>}

\subsubsection{Синтаксис}

\begin{lstlisting}
double pi = 3.14159;
int truncated = (int)pi;       // C-style cast
int also_works = int(pi);      // Функциональный стиль (то же самое)
\end{lstlisting}

\subsubsection{Как это работает внутри?}

C-style cast --- это \textbf{не отдельная операция}. Компилятор перебирает касты
в строгом порядке и применяет \textbf{первый подходящий}:

\begin{enumerate}
    \item \texttt{const\_cast} --- попробовать снять/добавить \texttt{const}
    \item \texttt{static\_cast} --- попробовать стандартное преобразование
    \item \texttt{static\_cast} + \texttt{const\_cast} --- комбинация
    \item \texttt{reinterpret\_cast} --- попробовать переинтерпретировать байты
    \item \texttt{reinterpret\_cast} + \texttt{const\_cast} --- самая опасная комбинация
\end{enumerate}

\subsubsection{Почему это плохо?}

\begin{lstlisting}
class Base { virtual ~Base() = default; };
class Derived : public Base { public: int x = 42; };

void example() {
    const Base* base_ptr = new Derived();

    // C-style cast: что именно здесь происходит?
    Derived* d = (Derived*)base_ptr;
    // Ответ: reinterpret_cast + const_cast (снял const И переинтерпретировал)
    // Вы это поняли, просто глядя на код? Вряд ли.

    // Явный C++ каст -- намерение прозрачно:
    const Derived* d2 = static_cast<const Derived*>(base_ptr);  // const сохранён
}
\end{lstlisting}

\subsubsection{Главная опасность}

C-style cast \textbf{никогда не выдаёт ошибку компиляции} для указателей. Он всегда
найдёт способ --- пусть даже через \texttt{reinterpret\_cast}:

\begin{lstlisting}
std::string str = "hello";
int* p = (int*)&str;   // Компилируется! Но это полная чушь.
                        // reinterpret_cast<int*>(&str) -- хотя бы видно, что опасно

// С C++ кастами:
// int* p = static_cast<int*>(&str);         // Ошибка компиляции! Несовместимы.
// int* p = reinterpret_cast<int*>(&str);    // Компилируется, но ЯВНО опасно
\end{lstlisting}

\textbf{Правило: в C++ коде C-style cast использовать нельзя. Точка.}

% -------------------------------------------------------
\subsection{\texttt{static\_cast} --- <<Безопасное преобразование>>}

\subsubsection{Суть}

\texttt{static\_cast} создаёт \textbf{новый объект} другого типа, используя
определённые правила преобразования. Это самый частый и самый безопасный каст.

Ключевое слово: \textbf{создаёт новый объект}. Старый остаётся нетронутым.

\textbf{Если вы не знаете, какой каст нужен --- скорее всего, нужен \texttt{static\_cast}.}

\subsubsection{Сценарий 1: Числовые преобразования}

\begin{lstlisting}
double pi = 3.14159;
int truncated = static_cast<int>(pi);  // 3 (дробная часть отброшена)

int big = 1'000'000;
short small = static_cast<short>(big);  // Возможна потеря данных (переполнение)

// Без каста компилятор выдаст предупреждение:
int x = 3.14;  // Warning: implicit conversion from 'double' to 'int'
int y = static_cast<int>(3.14);  // Предупреждения нет -- вы явно сказали "я знаю"
\end{lstlisting}

\textbf{Как \texttt{static\_cast<int>(3.14)} работает на уровне процессора?}

Это \textbf{не просто отбрасывание битов}. Числа \texttt{int} и \texttt{double}
хранятся совершенно по-разному:

\begin{lstlisting}[numbers=none]
double 3.14 в памяти (IEEE 754):
01000000 00001001 00011110 10111000 01010001 11101100 00000000 00000000
|знак   | экспонента |          мантисса                              |

int 3 в памяти:
00000000 00000000 00000000 00000011
\end{lstlisting}

\texttt{static\_cast} \textbf{вычисляет} новое значение: извлекает мантиссу
и экспоненту, считает целую часть и записывает в формате \texttt{int}. Это реальная
работа процессора (инструкция \texttt{cvttsd2si} на x86).

\subsubsection{Сценарий 2: Enum $\leftrightarrow$ int}

\begin{lstlisting}
enum class Color { Red = 0, Green = 1, Blue = 2 };

Color c = Color::Green;
int num = static_cast<int>(c);          // 1
Color back = static_cast<Color>(2);     // Color::Blue

// Без static_cast:
// int num = c;  // Ошибка! enum class не преобразуется неявно
\end{lstlisting}

\subsubsection{Сценарий 3: \texttt{void*} $\to$ конкретный тип}

\begin{lstlisting}
// Классический паттерн из C API (pthread, qsort и т.д.)
void callback(void* data) {
    // Мы ЗНАЕМ, что передали int*
    int* ptr = static_cast<int*>(data);
    std::cout << *ptr << "\n";
}

int main() {
    int value = 42;
    callback(&value);   // int* неявно -> void*
}
\end{lstlisting}

\subsubsection{Сценарий 4: Преобразование через конструктор}

\begin{lstlisting}
class Celsius {
public:
    double temp;
    explicit Celsius(double t) : temp(t) {}
};

double raw = 36.6;
// Celsius c = raw;  // Ошибка! Конструктор explicit
Celsius c = static_cast<Celsius>(raw);  // Явный вызов конструктора
\end{lstlisting}

\subsubsection{Сценарий 5: Приведение в иерархии наследования (downcast)}

\begin{lstlisting}
class Animal {
public:
    virtual ~Animal() = default;
    virtual void speak() = 0;
};

class Dog : public Animal {
public:
    void speak() override { std::cout << "Гав!\n"; }
    void fetch() { std::cout << "Принёс палку!\n"; }
};

void process(Animal* animal) {
    // Мы УВЕРЕНЫ, что это Dog (например, из контекста)
    Dog* dog = static_cast<Dog*>(animal);
    dog->fetch();  // Работает, если animal действительно Dog
                   // UB, если animal -- не Dog!
}
\end{lstlisting}

\textbf{Важно:} \texttt{static\_cast} \textbf{не проверяет} тип в рантайме.
Если вы ошиблись --- UB. Для безопасной проверки используйте \texttt{dynamic\_cast}.

\subsubsection{Чего \texttt{static\_cast} НЕ может}

\begin{lstlisting}
int x = 42;
// static_cast<double*>(&x);       // Ошибка! int* и double* несовместимы
// static_cast<std::string>(x);    // Ошибка! Нет конструктора string(int)

const int cx = 10;
// static_cast<int&>(cx);          // Ошибка! Не может снять const
\end{lstlisting}

% -------------------------------------------------------
\subsection{\texttt{reinterpret\_cast} --- <<Рентген памяти>>}

\subsubsection{Суть}

\texttt{reinterpret\_cast} \textbf{не создаёт новый объект}. Он говорит:
<<Возьми эти байты в памяти и считай, что это другой тип>>. Никаких вычислений,
никаких преобразований --- чистая переинтерпретация.

\subsubsection{Аналогия}

Представьте книгу на русском языке. \texttt{static\_cast} --- это \textbf{перевод}
на английский (создаётся новый текст с тем же смыслом). \texttt{reinterpret\_cast} ---
это попытка \textbf{прочитать русские буквы как английские} (ерунда, но технически
возможно).

\subsubsection{Пример 1: Заглядываем <<внутрь>> float}

\begin{lstlisting}
float f = 1.0f;

// Как float хранится в памяти?
// IEEE 754: 0 01111111 00000000000000000000000 = 0x3F800000

// НЕПРАВИЛЬНО -- UB по стандарту (нарушение strict aliasing):
int bits_bad = *reinterpret_cast<int*>(&f);

// ПРАВИЛЬНО -- через memcpy (стандарт разрешает):
int bits;
std::memcpy(&bits, &f, sizeof(float));
std::cout << std::hex << bits << "\n";  // 3f800000
\end{lstlisting}

\textbf{Strict aliasing rule:} обращение к объекту через указатель на <<чужой>> тип ---
UB (за исключением \texttt{char*}, \texttt{unsigned char*}, \texttt{std::byte*}).

\subsubsection{Пример 2: Указатель $\to$ число и обратно}

\begin{lstlisting}
int x = 42;
int* ptr = &x;

// Сохраняем адрес как число
uintptr_t addr = reinterpret_cast<uintptr_t>(ptr);
std::cout << "Адрес: 0x" << std::hex << addr << "\n";

// Восстанавливаем указатель из числа
int* restored = reinterpret_cast<int*>(addr);
std::cout << *restored << "\n";  // 42

// Это один из НЕМНОГИХ гарантированных сценариев для reinterpret_cast
\end{lstlisting}

\subsubsection{Пример 3: Работа с сетевыми пакетами}

\begin{lstlisting}
// Типичный паттерн при работе с сырыми данными из сети
struct PacketHeader {
    uint32_t magic;
    uint16_t version;
    uint16_t length;
};

void process_raw_data(char* buffer, size_t size) {
    // Технически UB из-за strict aliasing и alignment
    // Но широко используется на практике
    auto* header = reinterpret_cast<PacketHeader*>(buffer);

    // Безопасная альтернатива:
    PacketHeader header_safe;
    std::memcpy(&header_safe, buffer, sizeof(PacketHeader));

    if (header_safe.magic == 0xDEADBEEF) {
        // Обрабатываем пакет...
    }
}
\end{lstlisting}

\subsubsection{\texttt{reinterpret\_cast} НЕ может снять \texttt{const}}

\begin{lstlisting}
const int x = 42;
// int* p = reinterpret_cast<int*>(&x);  // Ошибка компиляции!
// Для снятия const нужен const_cast
\end{lstlisting}

\subsubsection{Когда \texttt{reinterpret\_cast} оправдан?}

Практически \textbf{никогда} в обычном коде. Легитимные случаи:
\begin{itemize}
    \item Преобразование указатель $\leftrightarrow$ \texttt{uintptr\_t}
    \item Преобразование между указателями на функции (с осторожностью)
    \item Низкоуровневый код (драйверы, аллокаторы, сериализация)
    \item Работа с C API, которое требует \texttt{void*}
\end{itemize}

% -------------------------------------------------------
\subsection{\texttt{const\_cast} --- <<Снятие защиты>>}

\subsubsection{Суть}

\texttt{const\_cast} --- единственный каст, который может \textbf{добавлять или
убирать} модификаторы \texttt{const} и \texttt{volatile}. Никакой другой каст
этого не умеет.

\subsubsection{Мягкая форма: добавление \texttt{const}}

Происходит \textbf{постоянно и неявно} --- вы даже не замечаете:

\begin{lstlisting}
std::vector<int> vec = {1, 2, 3};

// Когда вы вызываете size(), неявно происходит:
// vec (неконстантный) -> const vec& (константная ссылка внутри size())
size_t s = vec.size();  // size() объявлен как: size_t size() const;

// Это эквивалентно:
const std::vector<int>& const_ref = vec;  // Неявное добавление const
size_t s2 = const_ref.size();

// Явный const_cast (не нужен, но для понимания):
const std::vector<int>& ref = const_cast<const std::vector<int>&>(vec);
\end{lstlisting}

Добавление \texttt{const} --- \textbf{бесплатная} операция. Процессор не выполняет
никаких инструкций --- это чисто компиляторная проверка.

\subsubsection{Жёсткая форма: снятие \texttt{const}}

Вот здесь начинается опасность:

\begin{lstlisting}
void legacy_api(char* str) {
    // Старая C-функция, которая НЕ модифицирует строку,
    // но автор забыл написать const
    std::cout << str << "\n";
}

void modern_code(const std::string& s) {
    // Нужно вызвать legacy_api, но у нас const
    // legacy_api(s.c_str());  // Ошибка! c_str() возвращает const char*

    legacy_api(const_cast<char*>(s.c_str()));  // Снимаем const
    // Безопасно ТОЛЬКО если legacy_api реально не изменяет данные!
}
\end{lstlisting}

\subsubsection{Главное правило и UB}

\begin{lstlisting}
// БЕЗОПАСНО: объект изначально НЕ const, мы просто снимаем "лишний" const
int x = 42;
const int& ref = x;           // Добавили const
int& mutable_ref = const_cast<int&>(ref);  // Сняли const
mutable_ref = 100;             // OK -- x изначально не const
std::cout << x << "\n";       // 100

// UB: объект изначально ОБЪЯВЛЕН как const
const int y = 42;
int& bad_ref = const_cast<int&>(y);  // Компилируется...
bad_ref = 100;                        // UNDEFINED BEHAVIOR!
// Компилятор мог поместить y в read-only память
// Или заменить все обращения к y на литерал 42
\end{lstlisting}

\textbf{Почему UB?} Компилятор имеет право:
\begin{itemize}
    \item Поместить \texttt{const} объект в \textbf{read-only секцию} памяти
          (запись = segfault)
    \item \textbf{Заинлайнить} значение: заменить все \texttt{y} на литерал \texttt{42}
          при компиляции
    \item Оптимизировать код, \textbf{предполагая}, что значение никогда не меняется
\end{itemize}

\subsubsection{Реальный use case: избежание дублирования кода}

Классический паттерн для \texttt{operator[]}:

\begin{lstlisting}
class MyArray {
    int data[100];
public:
    // Константная версия -- основная логика
    const int& operator[](size_t idx) const {
        // Допустим, тут сложная проверка границ, логирование и т.д.
        if (idx >= 100) throw std::out_of_range("Index out of bounds");
        return data[idx];
    }

    // Неконстантная версия -- делегирует к константной
    int& operator[](size_t idx) {
        // Шаг 1: добавляем const к *this -> вызываем const-версию
        const auto& result = static_cast<const MyArray&>(*this)[idx];
        // Шаг 2: снимаем const с результата
        return const_cast<int&>(result);
        // Безопасно! Мы знаем, что *this не const (мы в неконстантном методе)
    }
};
\end{lstlisting}

Этот паттерн рекомендован Скоттом Мейерсом в <<Effective C++>>.

% -------------------------------------------------------
\subsection{\texttt{dynamic\_cast} --- <<Безопасный детектив>> (краткий обзор)}

\texttt{dynamic\_cast} проверяет тип \textbf{в рантайме} через таблицу виртуальных
функций (RTTI). Подробно изучается в теме ООП, но вот основная идея:

\begin{lstlisting}
class Base {
public:
    virtual ~Base() = default;  // Виртуальный деструктор ОБЯЗАТЕЛЕН
};

class Derived : public Base {
public:
    void special_method() { std::cout << "Я Derived!\n"; }
};

void safe_process(Base* base) {
    // dynamic_cast проверяет реальный тип в рантайме
    Derived* d = dynamic_cast<Derived*>(base);

    if (d != nullptr) {
        d->special_method();  // Гарантированно безопасно
    } else {
        std::cout << "Это не Derived\n";
    }
}

// Для ссылок: бросает std::bad_cast при неудаче
void safe_process_ref(Base& base) {
    try {
        Derived& d = dynamic_cast<Derived&>(base);
        d.special_method();
    } catch (const std::bad_cast& e) {
        std::cout << "Не Derived: " << e.what() << "\n";
    }
}
\end{lstlisting}

\begin{center}
\begin{tabular}{lcc}
\toprule
\textbf{Свойство} & \texttt{static\_cast} & \texttt{dynamic\_cast} \\
\midrule
Проверка в рантайме   & нет  & да \\
Нужен \texttt{virtual} & нет  & да \\
Скорость              & быстрый (0 затрат) & медленнее (RTTI) \\
При ошибке            & UB   & \texttt{nullptr} или исключение \\
\bottomrule
\end{tabular}
\end{center}

% -------------------------------------------------------
\subsection{Сводная таблица: какой каст когда}

\begin{center}
\begin{tabular}{lll}
\toprule
\textbf{Задача} & \textbf{Каст} & \textbf{Пример} \\
\midrule
\texttt{int} $\to$ \texttt{double}
    & \texttt{static\_cast}
    & \texttt{static\_cast<double>(42)} \\
\addlinespace
\texttt{enum} $\to$ \texttt{int}
    & \texttt{static\_cast}
    & \texttt{static\_cast<int>(Color::Red)} \\
\addlinespace
\texttt{Base*} $\to$ \texttt{Derived*} (уверен)
    & \texttt{static\_cast}
    & \texttt{static\_cast<Derived*>(base)} \\
\addlinespace
\texttt{Base*} $\to$ \texttt{Derived*} (не уверен)
    & \texttt{dynamic\_cast}
    & \texttt{dynamic\_cast<Derived*>(base)} \\
\addlinespace
Посмотреть байты объекта
    & \texttt{memcpy}
    & Через \texttt{char*} или \texttt{memcpy} \\
\addlinespace
Указатель $\leftrightarrow$ число
    & \texttt{reinterpret\_cast}
    & \texttt{reinterpret\_cast<uintptr\_t>(ptr)} \\
\addlinespace
Снять \texttt{const} для legacy API
    & \texttt{const\_cast}
    & \texttt{const\_cast<char*>(c\_str)} \\
\addlinespace
\texttt{void*} $\to$ конкретный тип
    & \texttt{static\_cast}
    & \texttt{static\_cast<int*>(void\_ptr)} \\
\addlinespace
<<Хочу чтобы скомпилировалось>>
    & \textbf{Ничего}
    & Подумайте ещё раз о дизайне \\
\bottomrule
\end{tabular}
\end{center}

% -------------------------------------------------------
\subsection{Упражнения}

\subsubsection{Упражнение 1: Найди ошибку}

Что не так в каждом фрагменте? Какой результат?

\begin{lstlisting}
// Фрагмент A
float f = 3.14f;
int* p = (int*)&f;
std::cout << *p << "\n";

// Фрагмент B
const std::string name = "Alice";
std::string& ref = const_cast<std::string&>(name);
ref = "Bob";

// Фрагмент C
int x = 65;
char* c = static_cast<char*>(&x);
\end{lstlisting}

\textbf{Решение:}

\begin{itemize}
    \item \textbf{Фрагмент A:} C-style cast выполняет
          \texttt{reinterpret\_cast<int*>(\&f)}. Чтение \texttt{*p} --- UB
          (нарушение strict aliasing: \texttt{int*} указывает на объект типа
          \texttt{float}). На практике выведет \texttt{1078523331}
          (\texttt{0x4048F5C3} --- IEEE~754 представление \texttt{3.14f}),
          но полагаться на это нельзя.

    \item \textbf{Фрагмент B:} UB! Объект \texttt{name} объявлен как \texttt{const}.
          Снятие \texttt{const} и модификация --- Undefined Behavior. Компилятор мог
          разместить \texttt{"Alice"} в read-only памяти (segfault). Или заоптимизировать
          и продолжать <<видеть>> \texttt{"Alice"}.

    \item \textbf{Фрагмент C:} Ошибка компиляции! \texttt{static\_cast} не может
          преобразовать \texttt{int*} $\to$ \texttt{char*} (несовместимые типы
          указателей). Для этого нужен \texttt{reinterpret\_cast<char*>(\&x)} ---
          и это один из немногих случаев, где такой каст легален (через \texttt{char*}
          можно читать байты любого объекта).
\end{itemize}

\subsubsection{Упражнение 2: Перепиши без C-style cast}

Замените C-style касты на правильные C++ касты:

\begin{lstlisting}
double d = 3.14;
int i = (int)d;

const char* s = "hello";
char* p = (char*)s;

Base* b = new Derived();
Derived* d = (Derived*)b;

void* raw = malloc(100);
int* arr = (int*)raw;
\end{lstlisting}

\textbf{Решение:}

\begin{lstlisting}
double d = 3.14;
int i = static_cast<int>(d);  // Числовое преобразование

const char* s = "hello";
char* p = const_cast<char*>(s);  // Снятие const (модификация = UB!)

Base* b = new Derived();
Derived* d = dynamic_cast<Derived*>(b);  // Безопасный downcast с проверкой
// Или static_cast<Derived*>(b) если вы на 100% уверены в типе

void* raw = malloc(100);
int* arr = static_cast<int*>(raw);  // void* -> конкретный тип
\end{lstlisting}

\subsubsection{Упражнение 3: Что выведет программа?}

\begin{lstlisting}
#include <iostream>

int main() {
    int x = 42;
    const int& cref = x;

    int& ref = const_cast<int&>(cref);
    ref = 100;

    std::cout << "x = " << x << "\n";
    std::cout << "cref = " << cref << "\n";
    std::cout << "ref = " << ref << "\n";

    return 0;
}
\end{lstlisting}

\textbf{Решение:}

Вывод:
\begin{lstlisting}[numbers=none]
x = 100
cref = 100
ref = 100
\end{lstlisting}

Это \textbf{не UB}, потому что \texttt{x} изначально объявлен как обычный
(не \texttt{const}) \texttt{int}. Мы лишь добавили \texttt{const} через ссылку,
а потом сняли его обратно. Все три имени (\texttt{x}, \texttt{cref}, \texttt{ref})
ссылаются на одну и ту же область памяти, поэтому изменение через \texttt{ref}
видно через все имена.

Было бы UB, если бы \texttt{x} был объявлен как \texttt{const int x = 42;}.

\subsubsection{Упражнение 4: Реализуй безопасное приведение}

Напишите шаблонную функцию \texttt{safe\_narrow}, которая преобразует числовой
тип с проверкой потери данных:

\begin{lstlisting}
template<typename To, typename From>
To safe_narrow(From value) {
    // Ваш код здесь
}

// Использование:
int big = 1'000'000;
short s = safe_narrow<short>(big);  // Должно бросить исключение!

double pi = 3.14;
int i = safe_narrow<int>(pi);  // Должно бросить исключение! (потеря .14)
\end{lstlisting}

\textbf{Решение:}

\begin{lstlisting}
#include <stdexcept>
#include <sstream>

template<typename To, typename From>
To safe_narrow(From value) {
    To converted = static_cast<To>(value);

    // Проверяем: если преобразовать обратно, получим ли исходное значение?
    if (static_cast<From>(converted) != value) {
        std::ostringstream oss;
        oss << "Narrowing conversion lost data: " << value
            << " became " << converted;
        throw std::runtime_error(oss.str());
    }

    return converted;
}

// Тесты:
int main() {
    // Безопасные преобразования
    short s1 = safe_narrow<short>(42);        // OK: 42 помещается в short
    int i1 = safe_narrow<int>(3.0);           // OK: 3.0 -> 3 без потерь

    // Опасные -- бросают исключение
    try {
        short s2 = safe_narrow<short>(1'000'000);  // Переполнение!
    } catch (const std::runtime_error& e) {
        std::cout << e.what() << "\n";
        // "Narrowing conversion lost data: 1000000 became 16960"
    }

    try {
        int i2 = safe_narrow<int>(3.14);  // Потеря дробной части!
    } catch (const std::runtime_error& e) {
        std::cout << e.what() << "\n";
        // "Narrowing conversion lost data: 3.14 became 3"
    }
}
\end{lstlisting}

Идея: \texttt{static\_cast} в обе стороны. Если после <<туда-обратно>> значение
изменилось --- данные потеряны.

\subsubsection{Упражнение 5: Детектив по кастам}

Определите, какой именно каст скрывается за каждым C-style cast. Помните порядок:
\texttt{const\_cast} $\to$ \texttt{static\_cast} $\to$ \texttt{static\_cast} +
\texttt{const\_cast} $\to$ \texttt{reinterpret\_cast} $\to$ \texttt{reinterpret\_cast} +
\texttt{const\_cast}.

\begin{lstlisting}
// 1.
double d = 3.14;
int i = (int)d;

// 2.
const int x = 42;
int* p = (int*)&x;

// 3.
int* ptr = (int*)0xDEADBEEF;

// 4.
class A {};
class B : public A {};
A* a = new B();
B* b = (B*)a;

// 5.
const double cd = 2.71;
int ci = (int)cd;
\end{lstlisting}

\textbf{Решение:}

\begin{enumerate}
    \item \texttt{static\_cast<int>(d)} --- стандартное числовое преобразование
          \texttt{double} $\to$ \texttt{int}. \texttt{const\_cast} не подходит
          (типы разные), \texttt{static\_cast} подходит первым.

    \item \texttt{const\_cast<int*>(\&x)} --- \texttt{\&x} имеет тип
          \texttt{const int*}, а цель \texttt{int*}. Разница только в \texttt{const},
          поэтому это \texttt{const\_cast}. Первый в очереди, и он подходит.

    \item \texttt{reinterpret\_cast<int*>(0xDEADBEEF)} --- целое число $\to$ указатель.
          \texttt{const\_cast} не подходит (разные типы), \texttt{static\_cast}
          не умеет \texttt{int} $\to$ указатель. Поэтому \texttt{reinterpret\_cast}.

    \item \texttt{static\_cast<B*>(a)} --- приведение вниз по иерархии (downcast).
          \texttt{const\_cast} не подходит, \texttt{static\_cast} подходит
          (есть отношение наследования).

    \item \texttt{static\_cast<int>(cd)} --- \texttt{cd} это значение (rvalue),
          \texttt{const} на значениях не мешает. Это просто \texttt{static\_cast<int>(cd)} ---
          числовое преобразование.
\end{enumerate}

\subsubsection{Упражнение 6: Реализуй type-punning безопасно}

Напишите функцию, которая показывает побитовое представление любого типа:

\begin{lstlisting}
template<typename T>
std::string bit_representation(const T& value);

// Использование:
std::cout << bit_representation(3.14f) << "\n";
// "01000000 01001000 11110101 11000011"

std::cout << bit_representation(42) << "\n";
// "00000000 00000000 00000000 00101010"
\end{lstlisting}

\textbf{Решение:}

\begin{lstlisting}
#include <cstring>
#include <string>
#include <bitset>
#include <sstream>

template<typename T>
std::string bit_representation(const T& value) {
    // Читаем байты через memcpy -- это БЕЗОПАСНО по стандарту
    unsigned char bytes[sizeof(T)];
    std::memcpy(bytes, &value, sizeof(T));

    std::ostringstream oss;
    for (size_t i = 0; i < sizeof(T); ++i) {
        if (i > 0) oss << " ";
        oss << std::bitset<8>(bytes[i]);
    }
    return oss.str();
}

int main() {
    std::cout << "float 3.14f:  " << bit_representation(3.14f) << "\n";
    std::cout << "int 42:       " << bit_representation(42) << "\n";
    std::cout << "double 3.14:  " << bit_representation(3.14) << "\n";
    std::cout << "char 'A':     " << bit_representation('A') << "\n";
}
\end{lstlisting}

Ключевой момент: мы используем \texttt{std::memcpy}, а \textbf{не}
\texttt{reinterpret\_cast}. Это единственный стандартно-безопасный способ
<<посмотреть на байты>> объекта. Чтение через \texttt{unsigned char*} /
\texttt{std::byte*} тоже допустимо по стандарту.

% -------------------------------------------------------
\subsection{Итоговая шпаргалка}

\begin{lstlisting}[numbers=none,basicstyle=\ttfamily\footnotesize]
+-----------------------------------------------------------+
|                  Какой каст выбрать?                      |
+-----------------------------------------------------------+
|                                                           |
|  Нужно преобразовать значение?                            |
|  (int->double, enum->int, void*->T*)                     |
|  -> static_cast                                           |
|                                                           |
|  Нужно снять/добавить const?                              |
|  -> const_cast                                            |
|                                                           |
|  Нужно посмотреть на байты / указатель<->число?           |
|  -> reinterpret_cast (а лучше memcpy)                     |
|                                                           |
|  Нужно безопасно привести тип в иерархии?                 |
|  -> dynamic_cast                                          |
|                                                           |
|  Не знаете какой?                                         |
|  -> static_cast (и подумайте, правильно ли спроектирован  |
|     код, если приходится кастовать)                        |
|                                                           |
|  C-style cast?                                            |
|  -> НИКОГДА в C++ коде                                    |
|                                                           |
+-----------------------------------------------------------+
\end{lstlisting}

% ===============================================================
\newpage
\section{Иерархия \texttt{std::exception}}

% -------------------------------------------------------
\subsection{Зачем нужна единая иерархия исключений?}

В разделе про \texttt{throw}/\texttt{catch} мы видели, что бросать можно
\textbf{любой} объект. Но если каждый программист будет бросать свои типы ---
наступит хаос: как ловить исключения из чужих библиотек?

Стандартная библиотека определяет \textbf{базовый класс} \texttt{std::exception}
и целое дерево наследников. Это позволяет:
\begin{itemize}
    \item Ловить \textbf{все} стандартные исключения одним \texttt{catch(const std::exception\&)}
    \item Ловить \textbf{конкретные} типы более специализированными блоками \texttt{catch}
    \item Создавать \textbf{свои} исключения, наследуясь от стандартных
\end{itemize}

% -------------------------------------------------------
\subsection{Класс \texttt{std::exception}}

Определён в заголовке \texttt{<exception>}.

\subsubsection{Интерфейс}

\begin{lstlisting}
class exception {
public:
    exception() noexcept;
    exception(const exception&) noexcept;
    exception& operator=(const exception&) noexcept;
    virtual ~exception() noexcept;

    virtual const char* what() const noexcept;
};
\end{lstlisting}

\begin{itemize}
    \item Все специальные методы помечены \texttt{noexcept} --- исключение
          \textbf{не должно} бросать исключений при своём создании/копировании.
    \item Единственный содержательный метод --- \texttt{what()}: возвращает
          C-строку с описанием ошибки.
    \item \texttt{what()} объявлен \texttt{virtual} --- наследники переопределяют его.
\end{itemize}

\subsubsection{Пример использования}

\begin{lstlisting}
#include <iostream>
#include <exception>

int main() {
    try {
        throw std::exception();  // Базовое исключение (редко бросают напрямую)
    } catch (const std::exception& e) {
        std::cout << e.what() << "\n";  // "std::exception" (зависит от реализации)
    }
}
\end{lstlisting}

% -------------------------------------------------------
\subsection{Дерево стандартных исключений}

\begin{lstlisting}[numbers=none,basicstyle=\ttfamily\footnotesize]
std::exception
|
+-- std::logic_error              (ошибки логики программы)
|   +-- std::invalid_argument     (неверный аргумент)
|   +-- std::domain_error         (значение вне области определения)
|   +-- std::length_error         (превышение максимального размера)
|   +-- std::out_of_range         (индекс за пределами)
|
+-- std::runtime_error            (ошибки времени выполнения)
|   +-- std::range_error          (результат вне диапазона)
|   +-- std::overflow_error       (арифметическое переполнение)
|   +-- std::underflow_error      (арифметическое антипереполнение)
|   +-- std::system_error         (ошибки ОС, C++11)
|
+-- std::bad_alloc                (нехватка памяти при new)
+-- std::bad_cast                 (неудачный dynamic_cast)
+-- std::bad_typeid               (typeid от nullptr)
+-- std::bad_exception            (нарушение спецификации исключений)
\end{lstlisting}

% -------------------------------------------------------
\subsection{\texttt{std::logic\_error} и \texttt{std::runtime\_error}}

Оба определены в \texttt{<stdexcept>} и хранят строку сообщения.

\subsubsection{Конструкторы}

\begin{lstlisting}
#include <stdexcept>

// Принимают std::string или const char*
std::logic_error   err1("ошибка логики");
std::runtime_error err2("ошибка выполнения");

// what() вернёт переданную строку
std::cout << err1.what() << "\n";  // "ошибка логики"
std::cout << err2.what() << "\n";  // "ошибка выполнения"
\end{lstlisting}

\subsubsection{Когда что бросать?}

\begin{center}
\begin{tabular}{lp{5cm}p{5cm}}
\toprule
\textbf{Тип} & \textbf{Смысл} & \textbf{Пример} \\
\midrule
\texttt{logic\_error}
    & Ошибка, которую \textit{можно было} предотвратить проверкой в коде
    & Передали отрицательный размер массива \\
\addlinespace
\texttt{runtime\_error}
    & Ошибка, которую \textit{нельзя} предсказать заранее
    & Файл не найден, сеть недоступна \\
\bottomrule
\end{tabular}
\end{center}

% -------------------------------------------------------
\subsection{Наследники \texttt{std::logic\_error}}

\subsubsection{\texttt{std::invalid\_argument}}

Аргумент функции имеет недопустимое значение.

\begin{lstlisting}
#include <stdexcept>
#include <string>

int parse_positive(const std::string& s) {
    int val = std::stoi(s);  // сам stoi может бросить invalid_argument
    if (val <= 0) {
        throw std::invalid_argument("expected positive number, got: " + s);
    }
    return val;
}

int main() {
    try {
        parse_positive("-5");
    } catch (const std::invalid_argument& e) {
        std::cout << "Ошибка: " << e.what() << "\n";
        // "Ошибка: expected positive number, got: -5"
    }
}
\end{lstlisting}

\subsubsection{\texttt{std::out\_of\_range}}

Индекс или ключ вне допустимого диапазона.

\begin{lstlisting}
#include <vector>
#include <stdexcept>

int main() {
    std::vector<int> v = {10, 20, 30};

    try {
        int val = v.at(100);  // бросает std::out_of_range
    } catch (const std::out_of_range& e) {
        std::cout << e.what() << "\n";  // сообщение зависит от реализации
    }

    // Для сравнения: v[100] -- UB, никакого исключения
}
\end{lstlisting}

\subsubsection{\texttt{std::length\_error}}

Попытка создать объект, превышающий максимально допустимый размер.

\begin{lstlisting}
#include <vector>
#include <stdexcept>

int main() {
    try {
        std::vector<int> v;
        v.reserve(v.max_size() + 1);  // бросает std::length_error
    } catch (const std::length_error& e) {
        std::cout << e.what() << "\n";
    }
}
\end{lstlisting}

\subsubsection{\texttt{std::domain\_error}}

Математическая функция получила значение вне области определения.

\begin{lstlisting}
#include <cmath>
#include <stdexcept>

double safe_sqrt(double x) {
    if (x < 0) {
        throw std::domain_error("sqrt of negative number");
    }
    return std::sqrt(x);
}
\end{lstlisting}

% -------------------------------------------------------
\subsection{Наследники \texttt{std::runtime\_error}}

\subsubsection{\texttt{std::overflow\_error} и \texttt{std::underflow\_error}}

\begin{lstlisting}
#include <stdexcept>
#include <limits>

unsigned int safe_multiply(unsigned int a, unsigned int b) {
    if (a != 0 && b > std::numeric_limits<unsigned int>::max() / a) {
        throw std::overflow_error("unsigned multiplication overflow");
    }
    return a * b;
}

int main() {
    try {
        auto result = safe_multiply(1'000'000'000u, 5u);
    } catch (const std::overflow_error& e) {
        std::cout << e.what() << "\n";
    }
}
\end{lstlisting}

\subsubsection{\texttt{std::system\_error} (C++11)}

Обёртка над ошибками операционной системы. Хранит \texttt{std::error\_code}.

\begin{lstlisting}
#include <system_error>
#include <fstream>

void open_file(const std::string& path) {
    std::ifstream f(path);
    if (!f.is_open()) {
        throw std::system_error(
            std::make_error_code(std::errc::no_such_file_or_directory),
            "cannot open: " + path
        );
    }
}

int main() {
    try {
        open_file("/nonexistent/file.txt");
    } catch (const std::system_error& e) {
        std::cout << "Код: " << e.code() << "\n";
        std::cout << "Что: " << e.what() << "\n";
    }
}
\end{lstlisting}

% -------------------------------------------------------
\subsection{Отдельные исключения (не наследуют \texttt{logic/runtime\_error})}

\subsubsection{\texttt{std::bad\_alloc}}

Бросается оператором \texttt{new} при нехватке памяти.

\begin{lstlisting}
#include <new>
#include <iostream>

int main() {
    try {
        // Попытка выделить абсурдно много памяти
        int* p = new int[1'000'000'000'000LL];
    } catch (const std::bad_alloc& e) {
        std::cout << "Нехватка памяти: " << e.what() << "\n";
    }

    // Альтернатива: new(std::nothrow) возвращает nullptr вместо исключения
    int* p = new(std::nothrow) int[1'000'000'000'000LL];
    if (p == nullptr) {
        std::cout << "Не удалось выделить память\n";
    }
}
\end{lstlisting}

\subsubsection{\texttt{std::bad\_cast}}

Бросается при неудачном \texttt{dynamic\_cast} к ссылке.

\begin{lstlisting}
#include <typeinfo>

class Base { public: virtual ~Base() = default; };
class Derived : public Base {};
class Other : public Base {};

int main() {
    Derived d;
    Base& ref = d;

    try {
        Other& o = dynamic_cast<Other&>(ref);  // Неудача!
    } catch (const std::bad_cast& e) {
        std::cout << e.what() << "\n";  // "std::bad_cast"
    }
}
\end{lstlisting}

% -------------------------------------------------------
\subsection{Создание собственных исключений}

\begin{lstlisting}
#include <stdexcept>
#include <string>

class HttpError : public std::runtime_error {
    int status_code_;
public:
    HttpError(int code, const std::string& msg)
        : std::runtime_error(msg), status_code_(code) {}

    int status_code() const noexcept { return status_code_; }
};

class NotFoundError : public HttpError {
public:
    explicit NotFoundError(const std::string& url)
        : HttpError(404, "Not found: " + url) {}
};

void fetch(const std::string& url) {
    throw NotFoundError(url);
}

int main() {
    try {
        fetch("/api/users/999");
    } catch (const HttpError& e) {
        std::cout << "HTTP " << e.status_code() << ": " << e.what() << "\n";
        // "HTTP 404: Not found: /api/users/999"
    }
    // Также можно поймать как std::runtime_error или std::exception
}
\end{lstlisting}

\textbf{Рекомендации:}
\begin{itemize}
    \item Наследуйтесь от \texttt{std::runtime\_error} или \texttt{std::logic\_error},
          а не от голого \texttt{std::exception} --- так вы получаете хранение строки
          бесплатно.
    \item Все методы исключений делайте \texttt{noexcept} (кроме конструктора,
          которому нужно скопировать строку).
\end{itemize}

% ===============================================================
\newpage
\section{\texttt{std::optional} (C++17)}

% -------------------------------------------------------
\subsection{Зачем нужен \texttt{std::optional}?}

Как обозначить <<значения нет>>? Исторически использовались костыли:

\begin{center}
\begin{tabular}{lp{4cm}p{5cm}}
\toprule
\textbf{Способ} & \textbf{Пример} & \textbf{Проблема} \\
\midrule
Магическое значение
    & \texttt{return -1;}
    & А если \texttt{-1} --- допустимый результат? \\
\addlinespace
Указатель
    & \texttt{return nullptr;}
    & Динамическая память, ownership неясен \\
\addlinespace
Пара \texttt{bool + value}
    & \texttt{std::pair<bool, T>}
    & Можно прочитать \texttt{value} при \texttt{false} \\
\addlinespace
Исключение
    & \texttt{throw NotFound\{\}}
    & Дорого, если <<нет значения>> --- обычная ситуация \\
\bottomrule
\end{tabular}
\end{center}

\texttt{std::optional<T>} --- это контейнер, который хранит \textbf{либо значение
типа \texttt{T}, либо ничего}. Определён в \texttt{<optional>}.

% -------------------------------------------------------
\subsection{Создание}

\subsubsection{Конструкторы}

\begin{lstlisting}
#include <optional>
#include <string>

// 1. Пустой optional
std::optional<int> empty;                    // нет значения
std::optional<int> also_empty = std::nullopt; // то же самое, явно

// 2. С значением
std::optional<int> opt1 = 42;               // содержит 42
std::optional<int> opt2(42);                // то же самое
std::optional<std::string> opt3("hello");   // содержит "hello"

// 3. in-place конструирование (без лишних копий)
std::optional<std::string> opt4(std::in_place, 5, 'x');  // содержит "xxxxx"
// Эквивалентно: std::string(5, 'x')

// 4. make_optional (выводит тип автоматически)
auto opt5 = std::make_optional(3.14);       // std::optional<double>
auto opt6 = std::make_optional<std::string>(5, 'x');  // "xxxxx"
\end{lstlisting}

% -------------------------------------------------------
\subsection{Проверка наличия значения}

\subsubsection{\texttt{has\_value()} и \texttt{operator bool()}}

\begin{lstlisting}
std::optional<int> opt = 42;
std::optional<int> empty;

// Способ 1: has_value()
if (opt.has_value()) {
    std::cout << "Есть значение\n";
}

// Способ 2: неявное преобразование в bool
if (opt) {
    std::cout << "Есть значение\n";   // то же самое
}

if (!empty) {
    std::cout << "Пусто\n";
}
\end{lstlisting}

Оба способа эквивалентны. \texttt{operator bool()} --- просто сокращение
для \texttt{has\_value()}.

% -------------------------------------------------------
\subsection{Доступ к значению}

\subsubsection{\texttt{operator*()} и \texttt{operator->()}}

Доступ \textbf{без проверки}. Если значения нет --- UB.

\begin{lstlisting}
std::optional<std::string> opt = "hello";

// operator* -- доступ к значению
std::string& s = *opt;
std::cout << s << "\n";       // "hello"

// operator-> -- доступ к членам
std::cout << opt->size() << "\n";  // 5

// UB! Нет значения, но обращаемся:
std::optional<int> empty;
// int x = *empty;  // UB -- поведение не определено
\end{lstlisting}

\subsubsection{\texttt{value()}}

Доступ \textbf{с проверкой}. Если значения нет --- бросает \texttt{std::bad\_optional\_access}.

\begin{lstlisting}
#include <optional>
#include <iostream>

int main() {
    std::optional<int> opt = 42;
    std::cout << opt.value() << "\n";  // 42

    std::optional<int> empty;
    try {
        int x = empty.value();  // Бросает исключение!
    } catch (const std::bad_optional_access& e) {
        std::cout << e.what() << "\n";  // "bad optional access"
    }
}
\end{lstlisting}

\subsubsection{\texttt{value\_or(default)}}

Возвращает значение, если есть, иначе --- переданное значение по умолчанию.

\begin{lstlisting}
std::optional<int> opt = 42;
std::optional<int> empty;

std::cout << opt.value_or(0) << "\n";    // 42 (есть значение)
std::cout << empty.value_or(0) << "\n";  // 0  (пусто, вернул default)

// Удобно для конфигурации:
std::optional<std::string> user_name;  // пользователь не указал имя
std::string name = user_name.value_or("Anonymous");
\end{lstlisting}

% -------------------------------------------------------
\subsection{Модификация}

\subsubsection{\texttt{emplace(args...)}}

Создаёт значение <<на месте>>, разрушая предыдущее (если было).

\begin{lstlisting}
std::optional<std::string> opt;

opt.emplace(5, 'a');         // Создаёт "aaaaa" внутри optional
std::cout << *opt << "\n";   // "aaaaa"

opt.emplace("hello");        // Старое разрушено, создано "hello"
std::cout << *opt << "\n";   // "hello"
\end{lstlisting}

\subsubsection{\texttt{reset()}}

Разрушает значение (если было), optional становится пустым.

\begin{lstlisting}
std::optional<std::string> opt = "hello";
std::cout << opt.has_value() << "\n";  // 1 (true)

opt.reset();
std::cout << opt.has_value() << "\n";  // 0 (false)

// Эквивалентно:
opt = std::nullopt;
\end{lstlisting}

\subsubsection{Присваивание}

\begin{lstlisting}
std::optional<int> opt;

opt = 42;              // Теперь содержит 42
opt = 100;             // Теперь содержит 100
opt = std::nullopt;    // Теперь пустой
opt = {};              // Тоже пустой
\end{lstlisting}

% -------------------------------------------------------
\subsection{Сравнение}

\texttt{std::optional} поддерживает все операторы сравнения:

\begin{lstlisting}
std::optional<int> a = 5;
std::optional<int> b = 10;
std::optional<int> c;       // пустой
std::optional<int> d;       // пустой

// optional с optional
std::cout << (a < b) << "\n";    // 1 (true): 5 < 10
std::cout << (a == b) << "\n";   // 0 (false)
std::cout << (c == d) << "\n";   // 1 (true): оба пустые
std::cout << (c < a) << "\n";    // 1 (true): пустой < любого значения

// optional с значением
std::cout << (a == 5) << "\n";   // 1 (true)
std::cout << (c == 5) << "\n";   // 0 (false): пустой != значение

// optional с nullopt
std::cout << (c == std::nullopt) << "\n";  // 1 (true)
std::cout << (a == std::nullopt) << "\n";  // 0 (false)
\end{lstlisting}

\textbf{Правила:}
\begin{itemize}
    \item Два пустых optional \textbf{равны}
    \item Пустой optional \textbf{меньше} любого непустого
    \item Два непустых сравниваются по значениям внутри
\end{itemize}

% -------------------------------------------------------
\subsection{Монадический интерфейс (C++23)}

C++23 добавил три метода для удобной цепочки операций.

\subsubsection{\texttt{and\_then(f)}}

Если есть значение --- применяет \texttt{f} (которая сама возвращает \texttt{optional}).
Если пусто --- возвращает пустой \texttt{optional}.

\begin{lstlisting}
// C++23
std::optional<std::string> get_env(const char* name);  // может вернуть nullopt
std::optional<int> parse_int(const std::string& s);    // может вернуть nullopt

// Цепочка: получить переменную среды -> распарсить как int
std::optional<int> port = get_env("PORT").and_then(parse_int);

// Без and_then (C++17):
std::optional<int> port2;
if (auto env = get_env("PORT")) {
    port2 = parse_int(*env);
}
\end{lstlisting}

\subsubsection{\texttt{transform(f)}}

Если есть значение --- применяет \texttt{f} и оборачивает результат в \texttt{optional}.
Если пусто --- возвращает пустой.

\begin{lstlisting}
// C++23
std::optional<std::string> name = "hello";
std::optional<size_t> len = name.transform([](const std::string& s) {
    return s.size();
});
// len содержит 5

std::optional<std::string> empty;
std::optional<size_t> no_len = empty.transform([](const std::string& s) {
    return s.size();
});
// no_len пуст
\end{lstlisting}

\subsubsection{\texttt{or\_else(f)}}

Если есть значение --- возвращает его. Если пусто --- вызывает \texttt{f}
(которая возвращает \texttt{optional}).

\begin{lstlisting}
// C++23
std::optional<std::string> primary_config();    // может быть nullopt
std::optional<std::string> fallback_config();   // может быть nullopt

// Попробовать primary, если нет -- fallback
auto config = primary_config().or_else(fallback_config);
\end{lstlisting}

% -------------------------------------------------------
\subsection{Где хранится значение?}

\texttt{std::optional<T>} хранит объект \textbf{внутри себя} (на стеке),
без динамической памяти. Упрощённая реализация:

\begin{lstlisting}
template<typename T>
class optional {
    alignas(T) unsigned char storage_[sizeof(T)];  // Сырая память
    bool has_value_;

    T& get() { return *reinterpret_cast<T*>(storage_); }
    const T& get() const { return *reinterpret_cast<const T*>(storage_); }

public:
    optional() : has_value_(false) {}

    optional(const T& value) : has_value_(true) {
        new (storage_) T(value);          // placement new
    }

    ~optional() {
        if (has_value_) get().~T();       // явный вызов деструктора
    }

    bool has_value() const { return has_value_; }
    T& operator*() { return get(); }
    T& value() {
        if (!has_value_) throw std::bad_optional_access();
        return get();
    }

    void reset() {
        if (has_value_) {
            get().~T();
            has_value_ = false;
        }
    }
};
\end{lstlisting}

\texttt{sizeof(std::optional<T>)} $\approx$ \texttt{sizeof(T)} + \texttt{sizeof(bool)}
+ выравнивание.

% -------------------------------------------------------
\subsection{Типичные паттерны использования}

\subsubsection{Функция, которая может не вернуть результат}

\begin{lstlisting}
#include <optional>
#include <vector>
#include <algorithm>

// Найти элемент в векторе
template<typename T>
std::optional<size_t> find_index(const std::vector<T>& v, const T& target) {
    for (size_t i = 0; i < v.size(); ++i) {
        if (v[i] == target) return i;
    }
    return std::nullopt;  // Не нашли
}

int main() {
    std::vector<int> v = {10, 20, 30, 40, 50};

    if (auto idx = find_index(v, 30)) {
        std::cout << "Найден на позиции " << *idx << "\n";  // 2
    } else {
        std::cout << "Не найден\n";
    }
}
\end{lstlisting}

\subsubsection{Ленивая инициализация (кэширование)}

\begin{lstlisting}
class Config {
    mutable std::optional<std::string> cached_name_;

public:
    const std::string& get_name() const {
        if (!cached_name_) {
            cached_name_ = expensive_computation();  // Считаем один раз
        }
        return *cached_name_;
    }

    void invalidate() { cached_name_.reset(); }  // Сбросить кэш
};
\end{lstlisting}

\subsubsection{Необязательные параметры структуры}

\begin{lstlisting}
struct UserProfile {
    std::string name;
    std::string email;
    std::optional<std::string> phone;    // Может отсутствовать
    std::optional<int> age;              // Может отсутствовать
};

void print_profile(const UserProfile& u) {
    std::cout << "Name: " << u.name << "\n";
    std::cout << "Email: " << u.email << "\n";

    if (u.phone) {
        std::cout << "Phone: " << *u.phone << "\n";
    }

    std::cout << "Age: " << u.age.value_or(-1) << "\n";
}
\end{lstlisting}

% ===============================================================
\newpage
\section{\texttt{std::expected} (C++23)}

% -------------------------------------------------------
\subsection{Зачем нужен \texttt{std::expected}?}

\texttt{std::optional} говорит: <<значение есть или его нет>>. Но \textbf{почему}
его нет? \texttt{optional} не даёт ответа.

\texttt{std::expected<T, E>} --- это контейнер, который хранит \textbf{либо
значение типа \texttt{T} (успех), либо ошибку типа \texttt{E} (неудача)}.
Определён в \texttt{<expected>} (C++23).

\begin{center}
\begin{tabular}{lp{4cm}p{4cm}}
\toprule
\textbf{Тип} & \textbf{Хранит} & \textbf{Когда использовать} \\
\midrule
\texttt{optional<T>}
    & \texttt{T} или ничего
    & <<Может не быть значения>> \\
\addlinespace
\texttt{expected<T,E>}
    & \texttt{T} или \texttt{E}
    & <<Может быть ошибка, и важно знать какая>> \\
\bottomrule
\end{tabular}
\end{center}

Это альтернатива исключениям для случаев, когда ошибка --- \textbf{ожидаемая}
часть логики (не исключительная ситуация):

\begin{center}
\begin{tabular}{lp{4.5cm}p{4.5cm}}
\toprule
& \textbf{Исключения} & \texttt{std::expected} \\
\midrule
Когда использовать
    & Действительно исключительные ситуации
    & Ожидаемые ошибки (парсинг, I/O, валидация) \\
\addlinespace
Стоимость ошибки
    & Дорого (раскрутка стека)
    & Дёшево (просто union) \\
\addlinespace
Видимость в API
    & Скрыта (надо читать документацию)
    & Явна в типе возвращаемого значения \\
\bottomrule
\end{tabular}
\end{center}

% -------------------------------------------------------
\subsection{Создание}

\subsubsection{Создание со значением (успех)}

\begin{lstlisting}
#include <expected>
#include <string>

// Неявно из значения
std::expected<int, std::string> result = 42;  // Содержит значение 42

// Явно
std::expected<int, std::string> result2(42);

// in-place конструирование
std::expected<std::string, int> result3(std::in_place, 5, 'x');  // "xxxxx"
\end{lstlisting}

\subsubsection{Создание с ошибкой (неудача)}

\begin{lstlisting}
// Через std::unexpected
std::expected<int, std::string> err = std::unexpected("file not found");

// Явное создание
std::expected<int, std::string> err2(std::unexpected("parse error"));

// in-place конструирование ошибки
std::expected<int, std::string> err3(std::unexpect, "timeout");
\end{lstlisting}

\textbf{Важно:} нельзя просто присвоить ошибку напрямую --- нужна обёртка
\texttt{std::unexpected}, чтобы компилятор отличил значение от ошибки:

\begin{lstlisting}
// Если T и E -- один тип, без unexpected была бы неоднозначность:
std::expected<std::string, std::string> result = "hello";  // Это значение!
std::expected<std::string, std::string> error =
    std::unexpected("oops");                               // Это ошибка!
\end{lstlisting}

% -------------------------------------------------------
\subsection{Проверка состояния}

\subsubsection{\texttt{has\_value()} и \texttt{operator bool()}}

\begin{lstlisting}
std::expected<int, std::string> ok = 42;
std::expected<int, std::string> err = std::unexpected("fail");

if (ok.has_value()) { /* true  */ }
if (ok)             { /* true  */ }

if (err.has_value()) { /* false */ }
if (!err)            { /* true  */ }
\end{lstlisting}

% -------------------------------------------------------
\subsection{Доступ к значению}

\subsubsection{\texttt{operator*()} и \texttt{operator->()}}

Доступ без проверки. Если содержит ошибку --- UB.

\begin{lstlisting}
std::expected<std::string, int> result = "hello";

std::string& s = *result;
std::cout << s << "\n";           // "hello"
std::cout << result->size() << "\n";  // 5

// UB:
std::expected<int, std::string> err = std::unexpected("fail");
// int x = *err;  // UB!
\end{lstlisting}

\subsubsection{\texttt{value()}}

Доступ с проверкой. Если ошибка --- бросает \texttt{std::bad\_expected\_access<E>}.

\begin{lstlisting}
std::expected<int, std::string> ok = 42;
std::cout << ok.value() << "\n";  // 42

std::expected<int, std::string> err = std::unexpected("file not found");
try {
    int x = err.value();  // Бросает исключение!
} catch (const std::bad_expected_access<std::string>& e) {
    // e.error() возвращает хранимую ошибку
    std::cout << "Ошибка: " << e.error() << "\n";  // "file not found"
}
\end{lstlisting}

\subsubsection{\texttt{value\_or(default)}}

\begin{lstlisting}
std::expected<int, std::string> ok = 42;
std::expected<int, std::string> err = std::unexpected("fail");

std::cout << ok.value_or(0) << "\n";   // 42
std::cout << err.value_or(0) << "\n";  // 0
\end{lstlisting}

% -------------------------------------------------------
\subsection{Доступ к ошибке}

\subsubsection{\texttt{error()}}

Возвращает хранимую ошибку. Если содержит значение (а не ошибку) --- UB.

\begin{lstlisting}
std::expected<int, std::string> err = std::unexpected("timeout");

std::cout << err.error() << "\n";  // "timeout"

// UB:
std::expected<int, std::string> ok = 42;
// auto e = ok.error();  // UB! Есть значение, а не ошибка
\end{lstlisting}

% -------------------------------------------------------
\subsection{Модификация}

\subsubsection{\texttt{emplace(args...)}}

\begin{lstlisting}
std::expected<std::string, int> result = std::unexpected(404);

// Заменяем ошибку на значение
result.emplace(5, 'x');              // Теперь содержит "xxxxx"
std::cout << *result << "\n";        // "xxxxx"
\end{lstlisting}

\subsubsection{Присваивание}

\begin{lstlisting}
std::expected<int, std::string> result = 42;

result = 100;                              // Значение 100
result = std::unexpected("new error");     // Теперь ошибка
result = 200;                              // Снова значение
\end{lstlisting}

% -------------------------------------------------------
\subsection{Монадический интерфейс (C++23)}

Те же три метода, что и у \texttt{optional}, плюс \texttt{or\_else} работает
с ошибкой.

\subsubsection{\texttt{and\_then(f)}}

Если значение --- применяет \texttt{f} (возвращающую \texttt{expected}).
Если ошибка --- пробрасывает ошибку.

\begin{lstlisting}
using Result = std::expected<int, std::string>;

Result parse_int(const std::string& s) {
    try {
        return std::stoi(s);
    } catch (...) {
        return std::unexpected("parse error: " + s);
    }
}

Result validate_positive(int x) {
    if (x > 0) return x;
    return std::unexpected("must be positive");
}

// Цепочка: парсим -> проверяем
Result result = parse_int("42").and_then(validate_positive);
// result содержит 42

Result fail = parse_int("abc").and_then(validate_positive);
// fail содержит ошибку "parse error: abc"
// validate_positive даже не вызывалась!
\end{lstlisting}

\subsubsection{\texttt{transform(f)}}

Если значение --- применяет \texttt{f}, оборачивает результат в \texttt{expected}.
Если ошибка --- пробрасывает.

\begin{lstlisting}
std::expected<int, std::string> result = 42;

auto doubled = result.transform([](int x) { return x * 2; });
// doubled содержит 84

std::expected<int, std::string> err = std::unexpected("fail");
auto still_err = err.transform([](int x) { return x * 2; });
// still_err содержит ошибку "fail"
\end{lstlisting}

\subsubsection{\texttt{or\_else(f)}}

Если значение --- возвращает его. Если ошибка --- вызывает \texttt{f} с ошибкой.

\begin{lstlisting}
using Result = std::expected<int, std::string>;

Result primary_source();
Result fallback_source(const std::string& error);

// Если primary не удался -- пробуем fallback, передавая ему текст ошибки
auto result = primary_source().or_else(fallback_source);
\end{lstlisting}

\subsubsection{\texttt{transform\_error(f)}}

Преобразует ошибку, не трогая значение. Полезно для конвертации типов ошибок.

\begin{lstlisting}
std::expected<int, int> result = std::unexpected(404);

// Преобразуем int-ошибку в строку
auto with_msg = result.transform_error([](int code) {
    return "HTTP error " + std::to_string(code);
});
// with_msg имеет тип std::expected<int, std::string>
// Содержит ошибку "HTTP error 404"
\end{lstlisting}

% -------------------------------------------------------
\subsection{Где хранится значение?}

Как и \texttt{optional}, хранит данные \textbf{внутри себя} без динамической
памяти. Внутри --- \texttt{union} из \texttt{T} и \texttt{E}:

\begin{lstlisting}
// Упрощённая реализация
template<typename T, typename E>
class expected {
    union {
        T value_;
        E error_;
    };
    bool has_value_;

public:
    expected(const T& val) : value_(val), has_value_(true) {}

    expected(std::unexpected<E> err)
        : error_(std::move(err).error()), has_value_(false) {}

    ~expected() {
        if (has_value_) value_.~T();
        else            error_.~E();
    }

    bool has_value() const { return has_value_; }
    T& operator*() { return value_; }
    E& error() { return error_; }

    T& value() {
        if (!has_value_) throw std::bad_expected_access(error_);
        return value_;
    }
};
\end{lstlisting}

\texttt{sizeof(std::expected<T, E>)} $\approx$
$\max(\texttt{sizeof(T)}, \texttt{sizeof(E)})$ + \texttt{sizeof(bool)} + выравнивание.

% -------------------------------------------------------
\subsection{\texttt{std::expected<void, E>}}

Специализация для функций, которые не возвращают значение, но могут вернуть ошибку.

\begin{lstlisting}
std::expected<void, std::string> save_file(const std::string& path) {
    std::ofstream f(path);
    if (!f.is_open()) {
        return std::unexpected("cannot open: " + path);
    }
    f << "data";
    if (f.fail()) {
        return std::unexpected("write failed");
    }
    return {};  // Успех (void)
}

int main() {
    auto result = save_file("/tmp/test.txt");
    if (!result) {
        std::cerr << "Ошибка: " << result.error() << "\n";
    }
}
\end{lstlisting}

% -------------------------------------------------------
\subsection{Типичные паттерны использования}

\subsubsection{Парсинг с информативной ошибкой}

\begin{lstlisting}
#include <expected>
#include <string>
#include <charconv>

enum class ParseError {
    empty_input,
    invalid_format,
    out_of_range,
};

std::expected<int, ParseError> parse_int(std::string_view sv) {
    if (sv.empty()) {
        return std::unexpected(ParseError::empty_input);
    }

    int result;
    auto [ptr, ec] = std::from_chars(sv.data(), sv.data() + sv.size(), result);

    if (ec == std::errc::invalid_argument) {
        return std::unexpected(ParseError::invalid_format);
    }
    if (ec == std::errc::result_out_of_range) {
        return std::unexpected(ParseError::out_of_range);
    }

    return result;
}

int main() {
    auto r1 = parse_int("42");
    if (r1) std::cout << "OK: " << *r1 << "\n";

    auto r2 = parse_int("abc");
    if (!r2) {
        switch (r2.error()) {
            case ParseError::empty_input:    std::cout << "Пусто\n"; break;
            case ParseError::invalid_format: std::cout << "Формат\n"; break;
            case ParseError::out_of_range:   std::cout << "Диапазон\n"; break;
        }
    }
}
\end{lstlisting}

\subsubsection{Цепочка операций}

\begin{lstlisting}
using Result = std::expected<int, std::string>;

Result read_config_value(const std::string& key);
Result validate_range(int value);
Result apply_multiplier(int value);

// Элегантная цепочка без вложенных if:
auto final_result = read_config_value("timeout")
    .and_then(validate_range)
    .and_then(apply_multiplier)
    .transform([](int x) { return x + 1; });

// Если любой шаг вернул ошибку -- она пробрасывается до конца
if (final_result) {
    std::cout << "Результат: " << *final_result << "\n";
} else {
    std::cerr << "Ошибка: " << final_result.error() << "\n";
}
\end{lstlisting}

% -------------------------------------------------------
\subsection{Сводная таблица: \texttt{optional} vs \texttt{expected} vs исключения}

\begin{center}
\begin{tabular}{lp{3.5cm}p{3.5cm}p{3.5cm}}
\toprule
& \texttt{optional<T>} & \texttt{expected<T,E>} & Исключения \\
\midrule
Информация об ошибке
    & Нет
    & Да (тип \texttt{E})
    & Да \\
\addlinespace
Стоимость ошибки
    & Нулевая
    & Нулевая
    & Высокая \\
\addlinespace
Стоимость успеха
    & Нулевая
    & Нулевая
    & Нулевая \\
\addlinespace
Можно забыть проверить
    & Да
    & Да
    & Нет \\
\addlinespace
Видимость в API
    & Явная
    & Явная
    & Скрытая \\
\addlinespace
Стандарт
    & C++17
    & C++23
    & C++98 \\
\bottomrule
\end{tabular}
\end{center}

% ===============================================================
\newpage
\section{CRTP --- Curiously Recurring Template Pattern}

% -------------------------------------------------------
\subsection{Что такое CRTP?}

\textbf{CRTP (Curiously Recurring Template Pattern)} --- паттерн, при котором
класс наследуется от шаблонного базового класса, \textbf{передавая себя}
в качестве шаблонного параметра:

\begin{lstlisting}
template<typename Derived>
class Base {
    // Base "знает" тип своего наследника
};

class MyClass : public Base<MyClass> {
    // MyClass наследуется от Base<MyClass>
};
\end{lstlisting}

Выглядит странно: класс наследуется от шаблона, параметризованного
\textit{самим собой}. Отсюда слово <<curiously>> (любопытно) в названии.

\subsubsection{Ключевая идея}

Базовый класс \texttt{Base<Derived>} может \textbf{вызывать методы наследника}
на этапе компиляции, без виртуальных функций:

\begin{lstlisting}
template<typename Derived>
class Base {
public:
    void interface() {
        // Приводим this к типу наследника и вызываем его метод
        static_cast<Derived*>(this)->implementation();
    }
};

class Concrete : public Base<Concrete> {
public:
    void implementation() {
        std::cout << "Concrete::implementation()\n";
    }
};

int main() {
    Concrete c;
    c.interface();  // "Concrete::implementation()"
}
\end{lstlisting}

\texttt{Base::interface()} вызывает \texttt{Derived::implementation()} через
\texttt{static\_cast}. Никаких виртуальных таблиц, никаких косвенных вызовов ---
всё разрешается \textbf{на этапе компиляции}.

% -------------------------------------------------------
\subsection{Зачем это нужно? Проблема виртуальных функций}

Виртуальные функции --- стандартный механизм полиморфизма в C++. Но у них
есть цена:

\begin{center}
\begin{tabular}{lp{5cm}p{5cm}}
\toprule
& \textbf{Виртуальные функции} & \textbf{CRTP} \\
\midrule
Полиморфизм
    & Динамический (runtime)
    & Статический (compile-time) \\
\addlinespace
Накладные расходы
    & \texttt{vptr} + косвенный вызов через vtable
    & Нулевые (всё инлайнится) \\
\addlinespace
Объект содержит \texttt{vptr}
    & Да (+8 байт на x86-64)
    & Нет \\
\addlinespace
Можно хранить в одном контейнере
    & Да (\texttt{vector<Base*>})
    & Нет (разные типы) \\
\addlinespace
Определяется в
    & Runtime
    & Compile-time \\
\bottomrule
\end{tabular}
\end{center}

CRTP --- это \textbf{статический полиморфизм}: поведение определяется на этапе
компиляции, а не в рантайме. Код получается таким же гибким, но \textbf{быстрее},
потому что компилятор видит конкретные типы и может инлайнить вызовы.

% -------------------------------------------------------
\subsection{Как работает \texttt{static\_cast} в CRTP?}

Центральный приём CRTP --- приведение \texttt{this} к типу наследника:

\begin{lstlisting}
template<typename Derived>
class Base {
public:
    Derived& self() {
        return static_cast<Derived&>(*this);
    }
    const Derived& self() const {
        return static_cast<const Derived&>(*this);
    }
};
\end{lstlisting}

\textbf{Почему это безопасно?} Потому что \texttt{Base<Derived>} используется
\textbf{только} как базовый класс для \texttt{Derived}. Объект всегда имеет тип
\texttt{Derived} (или его наследника), поэтому \texttt{static\_cast} вниз
по иерархии здесь корректен.

\textbf{Когда это опасно?} Если кто-то нарушит контракт:

\begin{lstlisting}
class A : public Base<A> { /* ... */ };
class B : public Base<A> { /* ... */ };  // BUG! B наследует Base<A>, а не Base<B>
                                          // static_cast<A*>(this) в объекте B = UB
\end{lstlisting}

Защита от этой ошибки рассмотрена в разделе~\ref{sec:crtp-protection}.

% -------------------------------------------------------
\subsection{Применение 1: Статический полиморфизм}

Классический пример --- интерфейс фигур без виртуальных функций.

\subsubsection{С виртуальными функциями (обычный подход)}

\begin{lstlisting}
class Shape {
public:
    virtual ~Shape() = default;
    virtual double area() const = 0;
    virtual void draw() const = 0;
};

class Circle : public Shape {
    double radius_;
public:
    explicit Circle(double r) : radius_(r) {}
    double area() const override { return 3.14159 * radius_ * radius_; }
    void draw() const override { std::cout << "Circle r=" << radius_ << "\n"; }
};
\end{lstlisting}

\subsubsection{С CRTP (статический полиморфизм)}

\begin{lstlisting}
template<typename Derived>
class Shape {
public:
    double area() const {
        return static_cast<const Derived*>(this)->area_impl();
    }
    void draw() const {
        static_cast<const Derived*>(this)->draw_impl();
    }
};

class Circle : public Shape<Circle> {
    double radius_;
public:
    explicit Circle(double r) : radius_(r) {}
    double area_impl() const { return 3.14159 * radius_ * radius_; }
    void draw_impl() const { std::cout << "Circle r=" << radius_ << "\n"; }
};

class Square : public Shape<Square> {
    double side_;
public:
    explicit Square(double s) : side_(s) {}
    double area_impl() const { return side_ * side_; }
    void draw_impl() const { std::cout << "Square s=" << side_ << "\n"; }
};
\end{lstlisting}

Использование с шаблонами вместо указателей на базовый класс:

\begin{lstlisting}
// Шаблонная функция -- принимает любую фигуру
template<typename ShapeT>
void print_area(const Shape<ShapeT>& shape) {
    std::cout << "Area: " << shape.area() << "\n";
}

int main() {
    Circle c(5.0);
    Square s(3.0);

    print_area(c);  // "Area: 78.5398"
    print_area(s);  // "Area: 9"

    c.draw();       // "Circle r=5"
    s.draw();       // "Square s=3"
}
\end{lstlisting}

\textbf{Что происходит при компиляции:}
\begin{itemize}
    \item \texttt{Shape<Circle>::area()} инлайнится: компилятор подставляет
          \texttt{Circle::area\_impl()} напрямую.
    \item Никакого \texttt{vptr}, никакого косвенного вызова через таблицу.
    \item \texttt{print\_area} инстанцируется для каждого типа фигуры ---
          получаются две отдельные, полностью оптимизированные функции.
\end{itemize}

% -------------------------------------------------------
\subsection{Применение 2: Добавление функциональности (mixins)}

CRTP позволяет базовому классу <<добавлять>> методы, которые используют
интерфейс наследника. Это называется \textbf{mixin}.

\subsubsection{Пример: автоматические операторы сравнения}

До C++20 (до \texttt{operator<=>}) приходилось писать все 6 операторов вручную.
CRTP позволяет вывести их из одного:

\begin{lstlisting}
// Mixin: требует от наследника только operator< и operator==
template<typename Derived>
class Comparable {
public:
    friend bool operator>(const Derived& a, const Derived& b) {
        return b < a;
    }
    friend bool operator<=(const Derived& a, const Derived& b) {
        return !(b < a);
    }
    friend bool operator>=(const Derived& a, const Derived& b) {
        return !(a < b);
    }
    friend bool operator!=(const Derived& a, const Derived& b) {
        return !(a == b);
    }
};

class Temperature : public Comparable<Temperature> {
    double value_;
public:
    explicit Temperature(double v) : value_(v) {}

    // Нужно определить только два оператора:
    friend bool operator<(const Temperature& a, const Temperature& b) {
        return a.value_ < b.value_;
    }
    friend bool operator==(const Temperature& a, const Temperature& b) {
        return a.value_ == b.value_;
    }
};

int main() {
    Temperature t1(36.6), t2(37.0);

    std::cout << (t1 < t2) << "\n";   // 1 -- определён вручную
    std::cout << (t1 > t2) << "\n";   // 0 -- сгенерирован CRTP
    std::cout << (t1 <= t2) << "\n";  // 1 -- сгенерирован CRTP
    std::cout << (t1 != t2) << "\n";  // 1 -- сгенерирован CRTP
}
\end{lstlisting}

\textbf{Примечание:} в C++20 для этого есть \texttt{operator<=>} (spaceship operator),
который решает ту же задачу средствами языка. Но паттерн mixin через CRTP применим
и к другим задачам.

\subsubsection{Пример: автоматический \texttt{operator!=} и постфиксный \texttt{++}}

\begin{lstlisting}
// Mixin для итераторов: из prefix ++ и == генерирует != и postfix ++
template<typename Derived>
class IteratorMixin {
public:
    // Постфиксный ++ через префиксный
    Derived operator++(int) {
        Derived old = static_cast<Derived&>(*this);
        ++static_cast<Derived&>(*this);  // Вызывает prefix ++
        return old;
    }

    friend bool operator!=(const Derived& a, const Derived& b) {
        return !(a == b);
    }
};

class Counter : public IteratorMixin<Counter> {
    int value_;
public:
    explicit Counter(int v) : value_(v) {}

    // Определяем только prefix ++ и ==
    Counter& operator++() { ++value_; return *this; }

    friend bool operator==(const Counter& a, const Counter& b) {
        return a.value_ == b.value_;
    }

    int get() const { return value_; }
};

int main() {
    Counter c(0);
    Counter end(5);

    while (c != end) {      // != сгенерирован CRTP
        std::cout << c.get() << " ";
        c++;                 // postfix ++ сгенерирован CRTP
    }
    // 0 1 2 3 4
}
\end{lstlisting}

% -------------------------------------------------------
\subsection{Применение 3: Подсчёт объектов}

CRTP позволяет создать счётчик экземпляров, \textbf{индивидуальный для каждого
класса}:

\begin{lstlisting}
template<typename Derived>
class Counter {
    static inline int count_ = 0;  // inline static (C++17)
public:
    Counter()  { ++count_; }
    Counter(const Counter&) { ++count_; }
    ~Counter() { --count_; }

    static int alive() { return count_; }
};

class Dog : public Counter<Dog> {
    std::string name_;
public:
    explicit Dog(std::string n) : name_(std::move(n)) {}
};

class Cat : public Counter<Cat> {
    std::string name_;
public:
    explicit Cat(std::string n) : name_(std::move(n)) {}
};

int main() {
    Dog d1("Шарик"), d2("Бобик");
    Cat c1("Мурка");

    std::cout << "Собак: " << Dog::alive() << "\n";  // 2
    std::cout << "Кошек: " << Cat::alive() << "\n";  // 1

    {
        Dog d3("Рекс");
        std::cout << "Собак: " << Dog::alive() << "\n";  // 3
    }
    std::cout << "Собак: " << Dog::alive() << "\n";  // 2
}
\end{lstlisting}

\textbf{Почему без CRTP это не работает?}

\begin{lstlisting}
// Без CRTP: один счётчик на ВСЕ классы
class CounterBad {
    static inline int count_ = 0;
public:
    CounterBad() { ++count_; }
    ~CounterBad() { --count_; }
    static int alive() { return count_; }
};

class Dog : public CounterBad {};
class Cat : public CounterBad {};

// Dog::alive() и Cat::alive() -- это ОДИН И ТОТ ЖЕ счётчик!
// Потому что Dog и Cat наследуют один и тот же CounterBad
\end{lstlisting}

С CRTP \texttt{Counter<Dog>} и \texttt{Counter<Cat>} --- это \textbf{разные классы},
каждый со своей статической переменной \texttt{count\_}.

% -------------------------------------------------------
\subsection{Применение 4: Singleton через CRTP}

\begin{lstlisting}
template<typename Derived>
class Singleton {
protected:
    Singleton() = default;
    ~Singleton() = default;

    Singleton(const Singleton&) = delete;
    Singleton& operator=(const Singleton&) = delete;

public:
    static Derived& instance() {
        static Derived inst;  // Meyer's singleton
        return inst;
    }
};

class Logger : public Singleton<Logger> {
    friend class Singleton<Logger>;  // Даём доступ к private конструктору
    Logger() = default;

public:
    void log(const std::string& msg) {
        std::cout << "[LOG] " << msg << "\n";
    }
};

class Config : public Singleton<Config> {
    friend class Singleton<Config>;
    Config() = default;

    int port_ = 8080;
public:
    int port() const { return port_; }
};

int main() {
    Logger::instance().log("started");     // "[LOG] started"
    std::cout << Config::instance().port() << "\n";  // 8080

    // Logger l;  // Ошибка компиляции! Конструктор protected
}
\end{lstlisting}

% -------------------------------------------------------
\subsection{Применение 5: Метод \texttt{clone()} без повторения кода}

При глубоком наследовании каждый класс должен переопределять \texttt{clone()}
с правильным типом. CRTP автоматизирует это:

\begin{lstlisting}
#include <memory>

class Animal {
public:
    virtual ~Animal() = default;
    virtual std::unique_ptr<Animal> clone() const = 0;
    virtual void speak() const = 0;
};

// CRTP-прослойка, которая реализует clone() за наследника
template<typename Derived>
class Cloneable : public Animal {
public:
    std::unique_ptr<Animal> clone() const override {
        // Derived ОБЯЗАН иметь копирующий конструктор
        return std::make_unique<Derived>(
            static_cast<const Derived&>(*this)
        );
    }
};

// Теперь наследникам не нужно писать clone() вручную:
class Dog : public Cloneable<Dog> {
    std::string name_;
public:
    explicit Dog(std::string n) : name_(std::move(n)) {}
    void speak() const override {
        std::cout << name_ << " говорит: Гав!\n";
    }
};

class Cat : public Cloneable<Cat> {
    std::string name_;
public:
    explicit Cat(std::string n) : name_(std::move(n)) {}
    void speak() const override {
        std::cout << name_ << " говорит: Мяу!\n";
    }
};

int main() {
    std::unique_ptr<Animal> dog = std::make_unique<Dog>("Шарик");
    dog->speak();  // "Шарик говорит: Гав!"

    auto cloned = dog->clone();
    cloned->speak();  // "Шарик говорит: Гав!" -- это КОПИЯ
}
\end{lstlisting}

Без CRTP пришлось бы копировать реализацию \texttt{clone()} в каждый класс.

% -------------------------------------------------------
\subsection{Применение 6: Статический интерфейс (<<концепты до концептов>>)}

До C++20 CRTP использовали для проверки интерфейса на этапе компиляции:

\begin{lstlisting}
template<typename Derived>
class Printable {
public:
    // Derived обязан иметь метод to_string()
    friend std::ostream& operator<<(std::ostream& os, const Derived& obj) {
        os << obj.to_string();
        return os;
    }
};

class Point : public Printable<Point> {
    double x_, y_;
public:
    Point(double x, double y) : x_(x), y_(y) {}

    std::string to_string() const {
        return "(" + std::to_string(x_) + ", " + std::to_string(y_) + ")";
    }
};

class Color : public Printable<Color> {
    int r_, g_, b_;
public:
    Color(int r, int g, int b) : r_(r), g_(g), b_(b) {}

    std::string to_string() const {
        return "rgb(" + std::to_string(r_) + "," +
               std::to_string(g_) + "," + std::to_string(b_) + ")";
    }
};

int main() {
    Point p(1.5, 2.3);
    Color c(255, 128, 0);

    std::cout << p << "\n";  // "(1.500000, 2.300000)"
    std::cout << c << "\n";  // "rgb(255,128,0)"
}
\end{lstlisting}

Наследник определяет \texttt{to\_string()}, а CRTP-база автоматически
генерирует \texttt{operator<<}. Если наследник забудет определить
\texttt{to\_string()} --- ошибка компиляции.

% -------------------------------------------------------
\subsection{Защита от ошибок}
\label{sec:crtp-protection}

\subsubsection{Проблема: неправильный шаблонный аргумент}

\begin{lstlisting}
class A : public Base<A> {};
class B : public Base<A> {};  // Ошибка! B передал A вместо B
                               // static_cast<A*>(this) на объекте B = UB
\end{lstlisting}

\subsubsection{Решение: проверка в конструкторе}

\begin{lstlisting}
template<typename Derived>
class Base {
    // Конструктор private, только Derived может вызвать (через наследование)
    Base() = default;
    friend Derived;  // Только Derived имеет доступ к конструктору
public:
    void interface() {
        static_cast<Derived*>(this)->implementation();
    }
};

class A : public Base<A> {};  // OK: A -- friend Base<A>
// class B : public Base<A> {}; // Ошибка компиляции!
                                // B не friend Base<A>, не может вызвать конструктор
\end{lstlisting}

\subsubsection{Решение C++20: \texttt{static\_assert} с концептами}

\begin{lstlisting}
template<typename Derived>
class Base {
public:
    Base() {
        // Проверяем на этапе компиляции, что Derived действительно наследник
        static_assert(std::is_base_of_v<Base<Derived>, Derived>,
                      "CRTP: Derived must inherit from Base<Derived>");
    }
};
\end{lstlisting}

% -------------------------------------------------------
\subsection{CRTP и множественное наследование}

Класс может наследовать \textbf{несколько} CRTP-баз --- каждая добавляет свою
функциональность:

\begin{lstlisting}
template<typename Derived>
class Serializable {
public:
    std::string serialize() const {
        return static_cast<const Derived*>(this)->serialize_impl();
    }
};

template<typename Derived>
class Loggable {
public:
    void log(const std::string& msg) const {
        std::cout << "[" << static_cast<const Derived*>(this)->class_name()
                  << "] " << msg << "\n";
    }
};

// User наследует оба mixin'а
class User : public Serializable<User>, public Loggable<User> {
    std::string name_;
    int age_;
public:
    User(std::string name, int age) : name_(std::move(name)), age_(age) {}

    // Для Serializable
    std::string serialize_impl() const {
        return "{\"name\":\"" + name_ + "\",\"age\":" + std::to_string(age_) + "}";
    }

    // Для Loggable
    const char* class_name() const { return "User"; }
};

int main() {
    User u("Alice", 25);
    std::cout << u.serialize() << "\n";  // {"name":"Alice","age":25}
    u.log("created");                     // [User] created
}
\end{lstlisting}

% -------------------------------------------------------
\subsection{CRTP vs виртуальные функции: когда что?}

\begin{center}
\begin{tabular}{lp{5.5cm}p{5.5cm}}
\toprule
\textbf{Критерий} & \textbf{Виртуальные функции} & \textbf{CRTP} \\
\midrule
Тип известен на этапе компиляции
    & Не обязательно
    & Да, всегда \\
\addlinespace
Нужен \texttt{vector<Base*>}
    & Да, это единственный способ
    & Нет (разные \texttt{Base<T>} --- разные типы) \\
\addlinespace
Критична производительность
    & Косвенный вызов, нет инлайнинга
    & Нулевые затраты, полный инлайнинг \\
\addlinespace
Количество типов растёт в рантайме (плагины)
    & Подходит
    & Не подходит \\
\addlinespace
Добавление общей функциональности (mixin)
    & Неудобно
    & Идеально \\
\addlinespace
Простота понимания
    & Привычно, интуитивно
    & Требует знания паттерна \\
\bottomrule
\end{tabular}
\end{center}

\textbf{Правило:}
\begin{itemize}
    \item Если набор типов известен \textbf{на этапе компиляции} и важна
          производительность --- CRTP.
    \item Если нужен \textbf{runtime-полиморфизм} (контейнер разнородных объектов,
          плагины, фабрики) --- виртуальные функции.
    \item Если нужно \textbf{добавить поведение} к разным классам (mixin) --- CRTP.
\end{itemize}

% -------------------------------------------------------
\subsection{Пример из реального мира: \texttt{std::enable\_shared\_from\_this}}

Стандартная библиотека использует CRTP в
\texttt{std::enable\_shared\_from\_this<T>}:

\begin{lstlisting}
#include <memory>

class Widget : public std::enable_shared_from_this<Widget> {
public:
    std::shared_ptr<Widget> get_ptr() {
        return shared_from_this();  // Метод из CRTP-базы
    }

    void do_work() {
        // Безопасно создаём shared_ptr на себя
        auto self = shared_from_this();
        // Можно передать self в асинхронную операцию,
        // гарантируя, что объект не будет уничтожен
    }
};

int main() {
    auto w = std::make_shared<Widget>();
    auto w2 = w->get_ptr();  // w и w2 разделяют владение

    std::cout << w.use_count() << "\n";  // 2
}
\end{lstlisting}

\texttt{enable\_shared\_from\_this<Widget>} --- это CRTP-база. Она хранит
внутри \texttt{weak\_ptr<Widget>}, который инициализируется при создании
\texttt{shared\_ptr<Widget>}. Метод \texttt{shared\_from\_this()} возвращает
\texttt{shared\_ptr}, созданный из этого \texttt{weak\_ptr}.

% -------------------------------------------------------
\subsection{Deducing this (C++23) --- замена CRTP?}

C++23 ввёл \textbf{deducing this} --- явный объектный параметр, который
упрощает многие сценарии, ранее требовавшие CRTP:

\begin{lstlisting}
// CRTP (до C++23):
template<typename Derived>
class Comparable {
    friend bool operator!=(const Derived& a, const Derived& b) {
        return !(a == b);
    }
};

class Int : public Comparable<Int> {
    int val_;
public:
    explicit Int(int v) : val_(v) {}
    friend bool operator==(const Int& a, const Int& b) {
        return a.val_ == b.val_;
    }
};
\end{lstlisting}

\begin{lstlisting}
// Deducing this (C++23):
class Comparable23 {
    // this передаётся явно, тип выводится автоматически
    bool operator!=(this const auto& self, const auto& other) {
        return !(self == other);
    }
};

class Int23 : public Comparable23 {
    int val_;
public:
    explicit Int23(int v) : val_(v) {}
    friend bool operator==(const Int23& a, const Int23& b) {
        return a.val_ == b.val_;
    }
};
\end{lstlisting}

C deducing this не нужен шаблонный параметр --- компилятор сам выводит
конкретный тип \texttt{self}. Код проще и безопаснее (нет риска передать
неправильный тип в шаблон базы).

\textbf{Однако CRTP не устарел:}
\begin{itemize}
    \item Deducing this доступен только с C++23
    \item CRTP работает и для статических членов (подсчёт объектов, singleton)
    \item Многие библиотеки и кодовые базы используют CRTP
    \item Понимание CRTP необходимо для чтения существующего кода
\end{itemize}

% -------------------------------------------------------
\subsection{Итого: шпаргалка по CRTP}

\begin{center}
\begin{tabular}{lp{9cm}}
\toprule
\textbf{Применение} & \textbf{Суть} \\
\midrule
Статический полиморфизм
    & Вызов методов наследника без \texttt{virtual} \\
\addlinespace
Mixin (добавление функциональности)
    & Генерация операторов, \texttt{operator<<}, \texttt{clone()} и~т.\,д. \\
\addlinespace
Подсчёт объектов
    & Отдельный \texttt{static} счётчик для каждого класса \\
\addlinespace
Singleton
    & Шаблонный \texttt{instance()} для любого класса \\
\addlinespace
\texttt{enable\_shared\_from\_this}
    & Получение \texttt{shared\_ptr} на \texttt{this} \\
\addlinespace
Избежание дублирования
    & Один \texttt{clone()}, один \texttt{serialize()} в базе \\
\bottomrule
\end{tabular}
\end{center}

\end{document}
